\documentclass[11pt]{article}
\usepackage[T1]{fontenc}
\usepackage{lmodern,amsmath,amssymb,amsthm,geometry,microtype,xcolor}
\usepackage[colorlinks=true,urlcolor=blue!55!black,linkcolor=black]{hyperref}
\geometry{margin=.93in}
\setlength{\emergencystretch}{2em}
\newtheorem{theorem}{Theorem}
\newtheorem{lemma}{Lemma}
\newcommand{\dd}{\,dl}
\newcommand{\dt}{\,dt}
\newcommand{\Rr}{\mathcal R}
\newcommand{\Ll}{\mathcal L}
\newcommand{\E}{\mathbb E}
\newcommand{\eps}{\varepsilon}
\newcommand{\norm}[1]{\lVert #1\rVert}
\newcommand{\repo}{https://github.com/TodManlaibaatar/OOD-reversals-theory/blob/737835735fa06a005eec743a9ec2493e247fe8bc}
\newcommand{\src}[2]{\href{\repo/09_canonical_regime_AN04/AN04_#1_v1.tex}{#2}}
\title{Initialized E5 on the closed WIT learned rectangle}
\author{AN04 \quad|\quad Version 1 for review}
\date{5 November 2026}
\begin{document}
\maketitle
\begin{abstract}
For each fixed critical initialization scale $c>0$ we prove the learned-endpoint depth comparison uniformly over WIT's closed rectangle. The finite-noise flow is centered at its own first measured half response. A uniform entrance proof, a common damping gap and a Gaussian lower gate estimate on a fixed-probability slab allow the signed donor history to be carried to the common witness time. The coherent comparison is the minimum for the specified parameters on WIT's common window.
\end{abstract}
\section{Claim, clocks and source results}\label{an04e5rv1:setup}
Use normalized time and the equally weighted mixture of $P_1=N(e_1,q^2I)$ and $P_2=N(\rho e_2,q^2I)$ with identity target $D$. Write $\E_j$ for expectation on $P_j$ and $\E$ for mixture expectation.
Let $(\rho,\zeta)$ range over the closed rectangle $\mathcal P_*$ in WIT \texttt{inc:AN06:v2:rectangle}. Set $\lambda=\rho^2$, $L=\lambda/2$, $\omega=(1-\lambda)/2$ and $\ell=\log(1/q)$. With $c>0$ fixed take $S_0=cq^{1/\omega}$. The initial vectors are $U_\alpha=W_\alpha=\sqrt{S_0}(\cos\alpha,\sin\alpha)$, with $dl=d\alpha/(2\pi)$. The quadrant and half-circle names always refer to this initial label. With $a_+=\max(a,0)$, put
\[
 \begin{gathered} f(D)=\int W(U^\top D)_+\dd,\qquad r=D-f,\qquad
 S=\int |U|^2\dd=\int |W|^2\dd,\\
 \Ll_j=\tfrac12\E_j|r|^2,\quad \Ll=(\Ll_1+\Ll_2)/2.\end{gathered}
\]
Balance holds label by label. For $\Theta=(U,W)$ the exact characteristics are
$U'=\mathsf R^\top W$, $W'=\mathsf R U$ with
$\mathsf R=\E[rD^\top{\bf1}_{U^\top D>0}]$ for the receiving label.
Thus $\Ll'=-\|\Theta'\|^2$ and $S'=2\E[f\cdot r]\le \E|D|^2/2$.
All unlabeled norms of parameters are in $L^2(dl)$. Set $Q_1=1+q^2$, $Q_2=Q=\lambda+q^2$ and
\[
 b_j=\E_j[D_j(U^\top D)_+]/Q_j,\quad v_j=\nabla_U b_j,\quad
 \theta_j=(v_j)_j,\qquad
 F=\begin{pmatrix}m&h\\z&n\end{pmatrix},\quad Fe_j=\int W b_j\dd.
\]
Let $t_h$ be the first half mass of right ancestry and $t_a=\inf\{t\ge t_h:n(t)=a\}$ for a weak response level $a$. Thus $t_{1/2}$ is measured from the actual flow. Use the central offset $s_*$ defining the common WIT interval $J$ and let $e_*=e_*(\rho,\zeta)=\min_{s\in J}\mathcal E_{\rm coh}(\rho,\zeta;s)$. Define $\xi=(\sin(q\zeta),-\cos(q\zeta))$ and $T=t_{1/2}+s_*$.
\begin{theorem}[Uniform rectangle endpoint]\label{an04e5rv1:main}
For every fixed $c>0$ there is a single $q_0(c)>0$ such that for every $0<q<q_0(c)$ and every $(\rho,\zeta)\in\mathcal P_*$, the first measured half response is finite, $t_{1/2}-t_h=\ell/L+O_c(1)$ uniformly, and at $T=t_{1/2}+s_*$,
\[
 m\ge\tfrac12,\qquad n\ge\tfrac14,\qquad
 \Ll_j(T)/\Ll_j(0)\le169/196\quad(j=1,2).
\]
The scaled depth satisfies
$\sup_{(\rho,\zeta)\in\mathcal P_*}e_q(T)\to-\infty$ as $q\downarrow0$, where $e_q(T)=(|\xi-f_T(\xi)|^2-1)/(2q^2)$.
In particular $e_q(T)<e_*$. The unscaled probe loss tends to $1/2$ uniformly.
\end{theorem}
\paragraph{Inputs from the theory repository.}
\src{trained_half_mass_bridge_donors}{HBD} and \src{bridge_return_before_weak_learning}{BRW}
give $t_h=\log(4/S_0)+o(1)$. At $t_h+b$ for each fixed $b\ge0$, the strong mass and regression tend to $M(b)=(1+e^{-b})^{-1}$, the right radial density tends to $4M(b)\cos^2\alpha$, and $4\Ll\to (1-M(b))^2+\lambda$. The Cartesian entrance envelopes have horizontal scale $H_0=\sqrt{S_0}e^{(1/2+q^2)t}$ and vertical scale $V_0=\sqrt{S_0}e^{(L+q^2)t}$. In particular both second components are $O_b(q)$ and the left radial mass per label is $O_b(q^2(1+t_h)^2+S_0)$ at fixed delay.
\src{core_output_descent}{COD} gives the resident rank ceiling and bottom-path bounds.
\src{transverse_balance_and_weak_clock}{TBW} gives the seed
\[
 q^{-2}\int_{\rm Q2}[(U_2+W_2)/2]^2\dd\ \longrightarrow\ c^{2\omega}(4e^b)^\lambda/8 .
\]
On compact Q2 interiors the seed divided by $q$ tends to $c^\omega(4e^b)^L\sin\alpha$ and both first components are $o_b(q)$. Fix $\eta$ as in TBW's small-hit proof, decreasing it if needed. On a fixed positive-length compact Q2 arc before $t_\eta$, $v=(U_2+W_2)/2>0$, $|U_1|,|W_1|,|U_2-W_2|/2\le v/2$, and $v'\ge\lambda v/4$ after a sufficiently late fixed entrance. Its loss budget gives $m>99/100$ there.
\src{residual_exchange_channels}{REC} and \src{cap_free_cross_response}{CFCR} give the Gaussian regression flow and shape estimates. We use the transverse energy identity of \src{gate_transport_and_transverse_memory}{GTM}. The \href{\repo/04_witnesses_persistence/AN06_learned_orbit_witnesses_v2.tex}{WIT v2} coherent orbit supplies the common interval and comparison. The proof follows the five lemmas of \href{\repo/09_canonical_regime_AN04/AN04_initialized_central_E5_v2.tex}{CE5 v2}, with the following parameter estimates and entrance verification.

\section{Scope and uniform entrance}\label{an04e5rv1:uniformsources}
Write $\mathcal P_*=\mathcal I_\rho\times\mathcal I_\zeta$ for WIT \texttt{rectangle}. For scalar envelopes we use the larger box with endpoints
\[
 \rho_-=13/20,\quad \rho_+=7/10,\quad z_-=87/50,\quad z_+=7/4,\qquad
 \lambda_-=169/400,\quad\lambda_+=49/100 .
\]
Thus $L_-=169/800\le L\le49/200=L_+$ and
$\omega_-=51/200\le\omega\le231/800=\omega_+$. Always retain the \emph{same} $\lambda$ and $L$ on the two sides of a cancellation. In estimates at fixed $c$ use $c_-=\min(c^{\omega_-},c^{\omega_+})$ and $c_+=\max(c^{\omega_-},c^{\omega_+})$. We may take $q^2\le1/100$, so $\lambda_-\le Q\le1/2$ and $1\le Q_1\le101/100$. All $O,o$ estimates below are uniform over $\mathcal P_*$, with previously chosen cutoffs fixed. A source theorem with a fixed family requires the following verification of this convention.

\paragraph{Domains of the imports.}
HBD \texttt{physical, clock, incoming, outer}, BRW \texttt{field, delay, budget, return} and COD \texttt{ceiling, exposure, field, trained} take each fixed $\rho\in[\rho_-,\rho_+]$, $c>0$, and $\zeta>a_\rho$. WIT \texttt{residentbounds} has $a_\rho<2/5< z_-$ throughout. BRW's protected \emph{first input moment} is in the proof of \texttt{return}, after \texttt{budget}; the return assertion itself is on fixed horizons. TBW \texttt{seed, smallhit, lossbudget, bulk} takes the critical initialized family on WIT's rectangle; its constants are stated for fixed family and witness. REC \texttt{shapes, four, features, grams}, CFCR \texttt{entry, flow, grams, scalar} and GTM \texttt{energy} have the same ratio box. For GTM we use its exact differentiated identity in the proof of \texttt{energy}; the REC/CFCR formulas used below follow at each state, independently of the stopping intervals in their other conclusions. WIT \texttt{centralmargins} is a central-orbit statement. Its replacement is the full-rectangle \texttt{uniformmargins}, used only after identifying the two clock phases in Section~\ref{an04e5rv1:guards}. No central initialized theorem is used as an off-central import.

\begin{lemma}[Fixed-delay entrance along varying parameters]\label{an04e5rv1:matching}
For fixed nominal caps and each fixed resident delay $b$, the physical bounds and compact-label entrance conclusions used below hold along any sequence $q_i\downarrow0$, $(\rho_i,\zeta_i)\to(\bar\rho,\bar\zeta)$ in the rectangle. Each extracted resident law solves the HBD--BRW field with the matched entrance at $\bar\rho$. Its tail, protected-moment and fixed-rank bounds can be chosen uniformly over all such laws and parameters.
\end{lemma}
\begin{proof}
Here are the ingredients of the extraction, including the early part whose length diverges. At nominal mass $A=S_0e^t$ in a fixed positive compact,
\[
 t=\ell/\omega+\log(A/c),\qquad
 e^{-\omega t}=q(c/A)^\omega,\qquad
 \sqrt{S_0}e^{Lt}=q\,c^\omega A^L .
\]
These are exact at each $i$; $e^{q^2t}=1+o(1)$ uniformly. HBD's component comparison with $H=\sqrt{S_0}e^{(1/2+q^2)t}$ has $\int H^2\le2A_{\max}$, $\sqrt{m_\alpha}\le H e^{H^2/2}$ and vertical envelope $C_{A_{\max}}H(e^{-\omega t}+q)$. The inhomogeneous convolution uses $1/\omega\le1/\omega_-$; Gaussian absolute moments in its fitted term are bounded by moments of $1+|Z|$ for standard $Z\in\mathbb R^2$ and $q\le1$. Thus the cap constants are independent of the ratio. The left squared radius per label is bounded by $C_{A_{\max}}(q^2(1+t)^2+S_0)$ at fixed nominal mass bounded below, as in HBD's left comparison: integrate the $O(qH)$ horizontal source to $O(qtH)$, the positive vertical component to $O(\sqrt{S_0}e^{(L+q^2)t}+qH/\omega_-)$, and propagate across a bounded nominal interval. The initial $\sqrt{S_0}$ term is $O(q)$ uniformly.

We detail why the incoming conditions travel with $\rho_i$. On compact Q1, at a fixed small cap $A_0$, CFS \texttt{an04cfsv1:laws} is uniform on the ratio box. Its target chart is $H_\rho^{-1}((c/A_0)^\omega\tan\alpha)$, and its target squared radius is asymptotic to $A_0\cos^2\alpha$. The chart depends continuously on $\rho$: in the defining quadrature $B_\rho'(w)\ge\omega_-$, the tail difference from $1/(\omega w)$ is bounded at large $w$ by $C K(-\rho_- w)/w^2$, and inversion on compact positive ranges is continuous. This tail bound also gives $H_\rho(w)/w\to1$ uniformly as $w\to\infty$. The trained chart error in CTE \texttt{an04ctev1:early, profile} at the cap is $C_I A_0(1+(c/A_0)^\omega)$ after scaling. Indeed the two comparison convolutions there use $1/\omega_-$ and $1/L_-$; the compact receiver has a common fixed positive distance from a vertical initial axis. The relative radial error is $e^{O(A_0)}$. On compact Q4 use RPS \texttt{an04rpsv1:q4}: arrival at the bounded scaled chart takes $2\ell+O_I(1)$, its attraction slope is at least $\omega_-$ and its remaining time is at least $(1/\omega_+-2)\ell-O(1)$. Both the reciprocal-speed bound on the outer segment and the root error $O(q^2)$ are uniform since $\rho\ge\rho_-$ and $1-\lambda\Phi\ge1-\lambda_+$. HBD's Q4 trained comparison uses the same $C_I A_0(1+(c/A_0)^\omega)$ and radial error. These errors tend to zero as $A_0\downarrow0$ since $\omega_+<1$.

At each fixed $A_0$, compact right radii stay bounded below on every later fixed nominal compact by the radial differential inequality. With the Cartesian envelope this gives bounded scaled charts, bounded chart velocities and the equicontinuity used in HBD. On compact receivers and donors the Gaussian kernels converge also as $\rho_i\to\bar\rho$: write the data as $\mu e_j+q Z$, use the HBD bounded-chart expansion with $\mu\in[\rho_-,1]$; the moving one-dimensional Gaussian gate thresholds converge, their boundary has probability zero and polynomial Gaussian moments dominate. Contributions of the omitted label strips are $O_{I,A_{\max}}(\text{strip length})$ in the scaled receiver equations by the same Cartesian envelope, with no dependence on $i$. Thus the label-path laws and separately tagged compact paths have subsequential limits on each fixed nominal window. Passing in their integral equations first on compacts and then removing strips gives HBD's field with coefficients $\bar\rho,\bar\lambda$. Take the small-cap matching just established at countably many fixed $A_0\downarrow0$. HBD's incoming conditions at the limit are $m/A\to\cos^2\alpha$, $A^{\bar\omega}(x,y)\to c^{\bar\omega}\tan\alpha(1,1)$ on Q1, and $(x,y)\to(a_{\bar\rho},a_{\bar\rho})$ on Q4. Its mass equation now gives $M=A/(4+A)$ and $m=4M\cos^2\alpha$. Strict crossing and the early spectral bound bracket the first half time uniformly. BRW extends the finite-horizon argument to the fixed cap $8e^b$. This proves fixed-delay transfer by contradiction and extraction; uniqueness of the full incoming angular law is unnecessary.

On Q2, CRA \texttt{an04crav1:outer} has exponent gap $1-\lambda/(1-\lambda)\ge2/51$. Thus at $A_0$ the target vertical seed in $q$ units equals $c^\omega A_0^L\sin\alpha(1+o(1))$ uniformly on each compact Q2 arc. The trained comparison has relative $O(A_0)$ error there. On $[T_{A_0},T_A]$ for fixed caps, the compact radius is $O(q)$, the off-diagonals are $O(q)$, the loss of either positive second component is at most $C_{A,A_0}q^2$ per unit time and its entry is $\ge c_{I,A_0}q$. This prevents contact with half the entry and makes the wrong gate exponentially small. TBW's sum equation is then $v'=Lv+O_{A,A_0}(q^2)$; integrating and then removing $A_0$ proves the stated seed. For the horizontal pair the small-cap contribution in $q$ units is $O(A_0^{1+L})$ and the continuation multiplier is at most $C_A(A/A_0)^{1/2+q^2}$; its limsup is $O_A(A_0^{1/2+L_-})$. Thus horizontal entry is $o(q)$. The limits are taken at fixed caps before $A_0\downarrow0$. On a compact Q3 arc both mean projections are strictly negative uniformly up to the early cap in the target comparison (Gaussian-tail drift integrated over $O(\ell)$); the fitted weight comparison and bounded later duration give $O_I(S_0)=o(q^2)$ as in COD \texttt{trained}. Fixed-delay Cartesian bounds pay for all endpoint strips.

Finally quantify the law-independent bounds in the resident regime. HBD's remote forcing budget is $C e^\tau(1+c^\omega e^{-\omega\tau})$ for $\tau\le\log8$, so the $J_*$ in \texttt{outer} is bounded by $C(8^{1+\omega_+}/(1+\omega_-)+8c_+)$. The lower right Q1 tail from $4^{-\omega}(c^\omega\tan\alpha-J_*)$ has normalized mass $\ge k(1+x)^{-3}$ above threshold $x$ in both coordinates, with a common $k>0$: integrate $4\cos^2\alpha\,dl$ on a vertical strip of width a fixed multiple of $(1+x)^{-1}$, using $c_->0$. The upper tail at half mass is $C(1+f^{-1/3})$ in physical rank $f$, by the uniform Cartesian bound. In BRW/COD the positive left eigenvector of the two-coordinate comparison stays bounded away from zero for $0\le f\le1/2$; its exponent is $\le L-\lambda_- f/16$. Integrating the damped ceiling against $df$ gives a common $B_0$ in $B_0e^{Lb}(1+b)^{-2/3}$ (the convolution denominator is bounded using $L_-$). In BRW's two-face proof take $b_1>\sqrt2 B_0$, $R_b=d+3b_1[(1+b)^{1/3}-1]$ with fixed $d>0$. The tail at $R_b$ times $M(b)\ge1/2$ times the lower input $d e^{Lb}/\sqrt2$ gives the common protected first moment $k'e^{Lb}/(1+b)$. The same estimate and ceiling give the fixed-band return with common time on a fixed compact, since $f\ge f_0>0$ there. COD's negative bottom input face $-H_1\sqrt{1+b}$ uses $\lambda_-H_1^2/2-L_+\ge1$ with $H_1$ then enlarged for the entry and prefactor. The negative bottom output is controlled by convolution of the preceding $B_0$ budget. For TBW's \emph{upper} bulk face one can take $F_b=C_F\log(1+b)/(1+b)$ with $\lambda_- C_F/32>3$ and log coefficient $B_F$ with $L_-B_F>2$, enlarging fixed constants if necessary. Indeed on $[b-2,b]$ the damped ceiling on $f\ge F_b$ is $o(e^{Lb}/(1+b))$; above $-\sqrt{b-B_F\log(1+b)}$ the Gaussian lower tail is $\gtrsim e^{-Lb}(1+b)^{LB_F-1/2}$. The downward speed on the negative strip dominates $\sqrt b$, and the positive part crosses zero in the preceding unit by the protected moment. Both faces are inward, including at contact. These comparisons use physical first moments. All bounds survive subsequences of parameters and compact tagged labels, proving the lemma.
\end{proof}

\section{Mass and time at the first small hit}\label{an04e5rv1:mass}
\paragraph{Choice of common majorants.}
For clarity, the following are bounds on the remaining parameter-sensitive comparisons in the proof. Generic constants obtained by adding, multiplying or integrating these majorants may depend on $c$, fixed arcs, the witness or an already chosen cutoff, in the stated order.
\begin{itemize}
\item On a fixed positive cap range $[a,A]$, bound $x^y$ by the largest of its four corner values for $x\in[a,A]$, $y\in[L_-,L_+]$ or $[\omega_-,\omega_+]$; its reciprocal uses the smallest. In particular $S_0/q^2\le c q^{1/\omega_+-2}$, and the scaled initial weighted output error is $O(\sqrt c\,q^{1/(2\omega_+)-1}\sqrt\ell)=o(1)$. The post-cap time is $\log(A/a)$, independent of $\rho$.
\item Any damped entrance integral dividing by $Q/2$, $a_1^*-a_2^*$, or $a_2^*$ uses $1/L_-$, $1/\omega_-$, or $1/L_-$ respectively. The strong deficit times the incoming growing envelope uses the strict rate gap in \eqref{an04e5rv1:damping}, cost at most $80000$ times the common entrance coefficient. Away from that entrance take the damping lower bound $(1/4)^2/(2S_*)$ for both $a_m,a_z$ with the mass bound $S_*$. In the rank ceiling the loss is at least $\lambda_-f/16$; a fixed compact rank has uniform gap $\nu_\delta=\lambda_-f_\delta/32$ to $L$. For negative resident output the discounted BRW integral split at $b/2$ costs $C L_-^{-1}(e^{-L_- b/2}+(1+b/2)^{-2/3})$. The charged-energy comparison retains $Q$ before exact clock cancellation; its subsequent large-delay loss uses $e^{-L_-b}$. For the incoming truncation $r(t_R)\le C e^{-L_-R}$; keep the sharper $Ce^{-LR}$ until multiplying by the centered propagator.
\item For a bounded centered horizon $|s|\le B_s$, a common Gronwall multiplier for the weak logistic comparison is $e^{C\lambda_+ B_s}$ with $C=1+2\sup|n|+2\sup|N|$. The uniform lower bound for the principal coefficient $w=Q(1-n)$ through $T$ is $\lambda_- \nu$ with $\nu$ from Section~\ref{an04e5rv1:guards}. For the negative-input excursion take $A_0=5/\rho_-$: the weak Gaussian tail charged by Cauchy--Schwarz and a fourth moment costs $Cq^{25/8}$, hence $O(q^2)$ per time. The strong gate costs $O(q)$ times the actual strong residual norm.
\item If $M$ bounds $\int |\mathsf R_{11}|$ on all prefixes and $a=2\sqrt{M(b_0)}$ denotes the horizontal entrance coefficient (here $M(b_0)$ is logistic), take $\epsilon_0=\min(1/2,ae^{-M}/8)$. The retained set constructed below admits e.g. $D=8(1+a e^M+e^M/a)$ for small $q$, common for all parameters and compact truncations. Fix $\kappa=\nu c_2/(4D)$ and the constants $d_g,\gamma$ of Lemma~\ref{an04e5rv1:gatelemma} before the truncation. The common factor on a centered interval of length $R+C_T$ is bounded by $2e^{L_+(R+|C_T|)}$. The incoming contribution after this amplification is bounded by $C_2e^{-L_-R}$ by combining exponents at the same $\rho$ first.
\end{itemize}

\paragraph{Uniform constants for the passage.}
All data and gate error constants in REC, CFCR and GTM may be taken on $q\le1/10$. Explicitly, a data moment of order $j$ is bounded by $\mathbb E(1+|Z|)^j$; inverse regression variances are at most $1/\lambda_-$; conditional wrong-principal-gate moments are bounded by Cauchy--Schwarz with these moments and $\exp[-\rho_-^2 (U_i/|U|)^2/(4q^2)]$. Bounds multiplied by $S$, $\sqrt S$ or $1+S$ retain these factors. Before removing the mass stop use $S\le4$. In the loss budget put $u=\lambda_-/400$. Choose $b_0$ so $(1-M(b_0))^2<u^2/8$ and $\eta<\min(1/8,u^2/(16\lambda_+))$; enlarge $b_0$ first as needed for the strong deficit. Lemma~\ref{an04e5rv1:matching} absorbs the fixed-delay $o(1)$ in this strict budget. Thus the three relevant residual norms are $\le u\le\lambda/400$; the full fourth norm is bounded by initial loss. Each off-diagonal and the $11$ field have size at most $\lambda_-/100\le\lambda/100$ under the loss estimates in the cone proof. As $(1-\eta)/2>7/16$, the error in $\mathsf R_{22}$ may be taken $<5\lambda_-/48$, giving $\mathsf R_{22}>\lambda/3$. The same inequalities cover both signs of each face and their intersections. The integrated energy and seed multipliers use $2/Q\le2/\lambda_-$ and $q(t_d-t_0)=o(1)$, with $d$ chosen small after $b_0$. On a fixed Q2 arc with $\sin\alpha\ge a_I>0$ the entrance amplitude is $\ge q c_-4^{L_-}a_I/2$ for $b_0\ge0$ and small $q$. The order-one lower amplitude at $t_d$ therefore has a common constant. Beyond $t_d$ the exponential-versus-linear comparison uses $\lambda_-/2$ in its squared growth and a fixed positive arc measure.

For the loss gap, write $\gamma_s=\sqrt{Q/Q_1}\in[3/5,3/4]$. The determinant identity below gives $2\gamma_s\det F\ge (3/5)(1/4)\eta-o(1)$ when $n\ge\eta/2$. In particular $\det F\ge\eta/20$ for small $q$. Let $B$ be the common mass bound up to the small hit, $M_A$ an upper bound for $\|FA\|_F$ from initial loss, and $a=\min(1,\eta/(2(1+\sup|n'|)))$, where $|n'|\le C(1+B)^2$ by the four-regression identity. Then $\sigma_{\min}(FA)\ge j=(\eta/20)\sqrt{\lambda_-}/M_A$. The leading lower bound in the output gradient is $j\sqrt{\lambda_-}(1-\eta)/(2\sqrt B)$. Taking half of this value, $g_0$, as a strict lower bound yields a common $\Delta=\min(\lambda_-/4,a g_0^2)>0$ in the gap. The one-sided excess over $Q$ at the start of this interval is uniformly $o(1)$: first fix an arbitrarily late delay, take the uniform entrance limit there and only then increase the delay. Afterwards $\sqrt Q-\sqrt{Q-\Delta}\ge\Delta/(2\sqrt{\lambda_++1/100})$. The pre-half error is at least $\sqrt{\lambda_-}/2$, so the dissipation over a fixed horizon dominates a fixed multiple of $\int_0^K(B+Cu)^{-1}\,du$ for sufficiently small $q$ at that fixed $K$. This fixes one $K$ and one clock offset bound for the rectangle. In particular
\[
 |2/Q-1/L|\,\ell\le 2q^2\ell/\lambda_-^2,\qquad
 e^{L(\ell/L)}=q^{-1},\qquad
 e^{(Q/2)(\ell/L)}=q^{-1}e^{q^2\ell/(2L)}\le 2/q .
\]
All final smallness thresholds are taken after the constants they enforce.

Write $t_0=t_h+b_0$ for a sufficiently late fixed entrance, $s=t-t_0$,
$\tau=t_\eta$, and $E=\tfrac12\int (U_2^2+W_2^2)\dd$. Stop at the first of
$\tau$, $S=4$, and $t_d=t_0+(2/Q)\log(d/q)$ for fixed $0<d<1$.
On this interval
\begin{equation}\label{an04e5rv1:energy}
 E'\le QE+C(q\sqrt E+q^2),\qquad |n|\le E+Cq\sqrt E .
\end{equation}
Indeed $b_2=(U_2)_++O(q|U|)$ gives the second bound by Cauchy--Schwarz.
The exact transverse identity is $E'=\E[D_2r^\top\partial_2 f+r_2 f_2]$.
On $P_1$, $\|\partial_2 f\|_\infty\le\sqrt{2SE}$, $D_2$ has norm $q$, and $r_2 f_2\le D_2^2/4$. On $P_2$, Euler's identity for the homogeneous map $f$ writes the integrand as
\[
 (D_1-f_1)(f_1-D_1\partial_1 f_1)+2r_2 f_2-D_1r_2\partial_1 f_2.
\]
The first product is at most $D_1^2(1-\partial_1 f_1)^2/4$ by completing the square. The last term costs $O(q\sqrt E)$, and the middle term has expectation at most $2Qn$. The cluster weight $1/2$ gives \eqref{an04e5rv1:energy}. Dividing by $2\sqrt{E+q^2}$ and integrating yields
\begin{equation}\label{an04e5rv1:envelope}
 \sqrt E\le C_{b_0}q e^{Qs/2},\quad
 \int_{t_0}^t\sqrt E\le C_{b_0}d,\quad \int_{t_0}^t E\le C_{b_0}d^2 .
\end{equation}
Projecting $S'$ on the principal cluster coordinates and completing the two residual squares gives $S'\le Q_1m(1-m)+Qn(1-n)+q^2/2$. Moreover
\[
 m'=a_m(1-m)+O(\sqrt E+q),\qquad
 a_m=\tfrac{Q_1}2\|\nabla m\|^2\ge Q_1m^2/S.
\]
Here $|\langle\nabla m,\nabla n\rangle|\le Cq S$ since $|v_1^\top v_2|\le Cq$, $|h|+|z|\le C(\sqrt E+q)$, and the projected shape error is $O(q)$ on the stop. Euler's identity $\langle\Theta,\nabla m\rangle=2m$ gives the lower bound. Integrating the damped equation bounds the positive deficit cost by $C(1-m(t_0))_++C\int(\sqrt E+q)$. Hence
\[
 S(t)\le S(t_0)+C(1-m(t_0))_++C_{b_0}(d+d^2)+o_q(1)<4 .
\]
Choose $b_0$ first using $S(t_0),m(t_0)\to M(b_0)$, then $d$, then $q$. This removes the mass stop.

\begin{lemma}[Q2 cone on both sides of the amplitude cutoff]\label{an04e5rv1:conelemma}
With $b_0$ sufficiently large, $\eta>0$ sufficiently small and fixed, and then $d>0$ fixed as above, the initialized first small hit satisfies
\begin{equation}\label{an04e5rv1:smallhit}
 \tau-t_h=\ell/L+O(1),\qquad
 \sup_{[t_0,\tau]}S+\int_{t_0}^{\tau}(|1-m|+|n|)\dt=O(1).
\end{equation}
On each fixed compact Q2 interior the positive cone below holds up to $\tau$ for small $q$.
\end{lemma}
\begin{proof}
Choose $b_0$ also large enough, and then $\eta<1/8$ small enough, that TBW \texttt{lossbudget} has square root less than $\lambda/400$ uniformly to the first $\eta$ hit. Fixed-delay BRW puts every earlier post-strong-half response up to $t_0$ below $\eta$ for small $q$. More explicitly, $4\Ll(t_0)\to(1-M(b_0))^2+\lambda$, and $\E_2 r_2^2\ge Q(1-n)^2\ge Q(1-\eta)^2$ on the no-hit interval. Thus
$\|r\|_{P_1},\|r_1\|_{P_2}\le\lambda/400$; $\|r_2\|_{P_2}\le C$ by initial loss. The regression identity then gives $m>99/100$. These choices can be made jointly with the preceding strict mass bound: its damping constants on $S\le4$, $m>99/100$ are absolute, while $d$ is chosen only after the fixed entrance.

Here is the budget for each field entry, including after $t_d$. Write $r_{ip}=\|r_i\|_{P_p}$ and $\Gamma={\bf1}_{U^\top D>0}$. The all-gate Cauchy--Schwarz bounds for the first three entries are
\[
 \begin{aligned}
 2|\mathsf R_{11}|&\le\sqrt{Q_1}\,r_{11}+q r_{12},\\
 2|\mathsf R_{12}|&\le q r_{11}+\sqrt Q\,r_{12},\\
 2|\mathsf R_{21}|&\le\sqrt{Q_1}\,r_{21}+q r_{22}.
 \end{aligned}
\]
Each is sufficient for $|\mathsf R_{ij}|\le\lambda/100$ for small $q$ (only $r_{22}$ uses the full loss bound). When $U_2/|U|\ge1/\sqrt2$, $\Pr_2(\Gamma=0)\le e^{-\rho^2/(4q^2)}$. The missing $P_2$ correlation costs at most $r_{22}(\E_2 D_2^4)^{1/4}\Pr_2(\Gamma=0)^{1/4}$. The ungated correlation is exactly $Q(1-n)$; the entire $P_1$ gate term costs at most $q r_{21}$. Hence
\[
 \mathsf R_{22}\ge Q(1-n)/2-Cq-Ce^{-c'/q^2}\ge\lambda/3
\]
on that cone through the full no-hit interval. These four estimates use loss only.

Let $v=(U_2+W_2)/2$, $d_v=(U_2-W_2)/2$. TBW's fixed-delay entrance on each chosen compact has $v/q$ bounded below by a positive constant, $U_1,W_1=o(q)$, $d_v=o(q)$. Thus $v>0$ and $|U_1|,|W_1|,|d_v|<v/2$ at $t_0$. On the closed faces $U_2,W_2\in[v/2,3v/2]$ and $U_2/|U|\ge1/\sqrt2$. The exact equations imply
\[
 v'=\mathsf R_{22}v+(\mathsf R_{12}W_1+\mathsf R_{21}U_1)/2
 \ge(\lambda/3-\lambda/200)v>\lambda v/4 .
\]
At $U_1=\pm v/2$, $D^+|U_1|\le |\mathsf R_{11}W_1|+|\mathsf R_{21}W_2|\le(2\lambda/100)v<v'/2$. At $W_1=\pm v/2$ use $|\mathsf R_{11}U_1|+|\mathsf R_{12}U_2|$ for exactly the same bound. At $d_v=\pm v/2$, its derivative is $-\mathsf R_{22}d_v+(\mathsf R_{12}W_1-\mathsf R_{21}U_1)/2$, so $D^+|d_v|\le(-\lambda/6+\lambda/200)v<v'/2$. These strict inequalities hold at intersections of any of the six faces. $v\ge v(t_0)e^{\lambda(t-t_0)/4}$ excludes the zero face. Each label therefore stays in the cone until the first hit, inclusively by continuity. In particular an arrival at $t_d$ with $\tau>t_d$ is an interior continuation time for the same loss argument; a simultaneous hit ends it there.

To obtain the sharper clock use the bounded-mass estimates \emph{only} up to $\min(t_d,\tau)$. REC's shape decomposition in the receiving field gives
$\mathsf R_{12}=\tfrac12Q_1(1-m)(v_1)_2-\tfrac12 Q h(v_2)_2+O(q)$ and
$\mathsf R_{21}=-\tfrac12Q_1 z(v_1)_1+\tfrac12Q(1-n)(v_2)_1+O(q)$.
Each remainder is bounded by $\frac12\sum_p\|e^{(p)}\|_{P_p}\|D_j\|_{P_p}=O(q)$ on $S\le4$. Thus $|\mathsf R_{12}|+|\mathsf R_{21}|\le C(\sqrt E+q)$. Combining the cone faces, the sharper lower bound $Q(1-n)/2-Cq$ and \eqref{an04e5rv1:energy} yields
$v'\ge[Q/2-C(E+\sqrt E+q)]v$ on this restricted interval. Integration using \eqref{an04e5rv1:envelope} gives a lower multiplier $c_{b_0,d}e^{Q(t-t_0)/2}$ since $q(t_d-t_0)=o(1)$. At $t_d$ this exponential equals $d/q$. Therefore, if $\tau>t_d$, a \emph{fixed positive-measure} interior arc $J_0$ has $v(t_d)\ge k_{b_0,d,J_0}>0$. For $0\le u\le\tau-t_d$ the already checked loss-budget cone and $v^2\le |U|^2$ by balance give
\[
 |J_0|k_{b_0,d,J_0}^2 e^{\lambda u/2}
 \le\int_{J_0}v(t_d+u)^2\dd\le S(t_d+u)\le4+Cu .
\]
Here $|J_0|=\int_{J_0}dl$, and the last inequality uses the global $S'$ bound. The exponential exceeds the line at a fixed finite $u$, so the hit occurs within $O(1)$ after $t_d$. This argument needs no sharpened off-diagonal estimate beyond $t_d$.

If $\tau\le t_d$, \eqref{an04e5rv1:energy} and $n(\tau)=\eta$ give $\sqrt{E(\tau)}\ge c_\eta>0$ for small $q$, so \eqref{an04e5rv1:envelope} gives $\tau-t_0\ge(2/Q)\ell-O(1)$. The upper bound is $t_d$ in this case. If $\tau>t_d$, both bounds follow from $t_d-t_0=(2/Q)\log(d/q)$ and the bounded continuation. Because $q^2\ell=o(1)$, $(2/Q)\ell=\ell/L+o(1)$. Up to $\min(\tau,t_d)$, the damped $m$ equation integrated in absolute value gives $\int |1-m|\le C(|1-m(t_0)|+\int(\sqrt E+q))$, and \eqref{an04e5rv1:energy} bounds $\int |n|$. The possible remaining interval is bounded and both regressions are bounded by loss. This proves \eqref{an04e5rv1:smallhit}. A fixed further continuation has bounded mass by the same linear inequality.
\end{proof}


\section{A fixed loss gap and the first half response}\label{an04e5rv1:half}
Let $A=\operatorname{diag}(\sqrt{Q_1},\sqrt Q)$, $u=A(b_1,b_2)^\top$, and $D_{\Ll}=4\Ll$.
Cluster projection gives $D_{\Ll}=\|(I-F)A\|_F^2+\chi$ with $0\le\chi\le Cq^2(1+S)^2$. Indeed $r$ is $(1+S)$-Lipschitz on each cluster, and approximation by its mean-center value times $D_j/\E_jD_j$ leaves $O(q(1+S))$ in $L^2(P_j)$. Since $FA=\int Wu^\top\dd$, $\int uu^\top\dd\succeq (FA)^\top(FA)/S$. The output gradient is $\dot W=(I-F)Au/2+N_W$ with $\|N_W\|\le Cq\sqrt S(1+S)$. Consequently
\begin{equation}\label{an04e5rv1:gram}
 \|\dot W\|\ge\frac{\sigma_{\min}(FA)}{2\sqrt S}\|(I-F)A\|_F-Cq\sqrt S(1+S).
\end{equation}
The bounded mass in \eqref{an04e5rv1:smallhit} bounds $|n'|$ uniformly. There is a fixed $a>0$ with $\eta/2\le n\le\eta$ on $[\tau-a,\tau]$. For every sufficiently late fixed $b$ the entrance at $t_h+b$ precedes this interval. Loss monotonicity and then $b\to\infty$ give $\limsup_{q\to0}(D_{\Ll}(\tau-a)-Q)\le0$. With $\gamma=\sqrt{Q/Q_1}$, the exact square completion reads
\[
 2\gamma\det F=(Q-D_{\Ll}+\chi)/Q_1+
 (m-1+\gamma n)^2+(z-\gamma h)^2+2\gamma(1-\gamma)n .
\]
Thus $\det F$ is bounded below on a fixed time interval. Since $\|FA\|$ is bounded there, its least singular value is bounded below too. The error norm in \eqref{an04e5rv1:gram} is at least $\sqrt Q(1-\eta)$, so $D_{\Ll}'\le-c_0<0$ there. In particular
\begin{equation}\label{an04e5rv1:gap}
 D_{\Ll}(\tau)\le Q-\Delta
\end{equation}
for a fixed $\Delta>0$. Afterwards $\sigma_{\min}(FA)\ge\sqrt Q-\sqrt{Q-\Delta}$ by the triangle inequality. Before half response the remaining error norm is at least $\sqrt Q/2$. On $[\tau,\tau+K]$, $S\le B+C(t-\tau)$. Choose $K$ so that the dissipation lower bound proportional to $\int_0^K (B+Cu)^{-1}du$ exceeds all available loss, then take $q$ small enough to absorb the error in \eqref{an04e5rv1:gram} throughout that fixed interval. This forces
\begin{equation}\label{an04e5rv1:halfclock}
 0\le t_{1/2}-\tau\le K,\qquad t_{1/2}=t_h+\ell/L+O(1).
\end{equation}
The argument requires neither monotonicity of $n$ nor a small cross-response cap.

\section{Axis separation and the learning guards}\label{an04e5rv1:guards}
\begin{lemma}[Charged energy and iterated axis separation]\label{an04e5rv1:axislemma}
Fix a finite $K$, $t_f=t_{1/2}+K$, and a closed interior arc $J_2\Subset{\rm Q2}$.
Charge coordinate $i=1$ on $D_1=J_2$ and coordinate $i=2$ on $D_2=J_2^c$ in both layers. In this lemma $D_i$ denotes a label set when used as an integration domain. Let $P$ be this projection and $E_J=\|P\Theta\|^2/2$. There is a constant independent of $J_2$ such that on $[t_0,t_f]$
\begin{equation}\label{an04e5rv1:defect}
 E_J'\le [Q+C(|n|+|1-m|)]E_J+C(q\sqrt{E_J}+q^2).
\end{equation}
The supremum of $\sqrt{E_J}$ on $[t_h+b,t_f]$ has iterated upper limit zero on taking first $q\downarrow0$, then $b\to\infty$, then $J_2\uparrow{\rm Q2}$. The right vertical norm, $|z|+|h|$ and $|T_1|$ tend uniformly to zero on $[t_0,t_f]$, where $T_1=\int W_1W_2|v_1|^2\dd$.
\end{lemma}
\begin{proof}
Write $g_{ij}=\nabla F_{ij}=(W_i v_j,e_i b_j)$ and $j'=3-j$. The small-hit clock and fixed loss gap in Sections~\ref{an04e5rv1:mass}--\ref{an04e5rv1:half} give a common $S\le S_*$ and a common bound for $\int_{t_0}^{t_f}(|n|+|1-m|)$. Loss monotonicity bounds each regression uniformly. Gaussian homogeneity gives the exact identity $b_j=U_j(v_j)_j+U_{j'}(v_j)_{j'}$, also at $U=0$ with the zero-gate convention. REC gives $|(v_j)_{j'}|\le Cq$ and $|b_j-(U_j)_+|\le Cq|U|$. These identities retain the derivative even on a center-gate boundary.

For $i=j$, the pairing with $P\Theta$ on $D_j$ has integrand
$W_j(b_j+U_j(v_j)_j)$; on $D_{j'}$ it is $W_j U_{j'}(v_j)_{j'}$. Hence
\[
 \langle P\Theta,g_{jj}\rangle
 =2\int_{D_j}W_j(U_j)_+\dd+e_j,\qquad
 |e_j|\le Cq\left(\int_{D_j}|W_j||U|\dd+
              \int_{D_{j'}}|W_j||U_{j'}|\dd\right)\le Cq\sqrt{E_J}.
\]
The last estimate uses $S_*$ and, in each integral, one charged factor. The leading integral in absolute value is bounded by $\int_{D_j}(W_j^2+U_j^2)\dd\le2E_J$ including its coefficient two. This accounts separately for $g_{11}=g_m$ and $g_{22}=g_n$.

For $i\ne j$, on $D_i$ the integrand is $W_i(b_j+U_i(v_j)_i)$; on $D_j$ it is $W_iU_j(v_j)_j$. Subtract $F_{ij}$ and use Euler's identity on $D_j$ to get
\[
 \langle P\Theta,g_{ij}\rangle-F_{ij}
 =\int_{D_i}W_iU_i(v_j)_i\dd-\int_{D_j}W_iU_i(v_j)_i\dd,\qquad
 |\langle P\Theta,g_{ij}\rangle-F_{ij}|\le CqS.
\]
Moreover $|F_{ij}|\le C\sqrt{E_J}+Cq$: approximate $b_j$ by $(U_j)_+$ and use $W_i$ on $D_i$, $U_j$ on $D_j$ as the charged factor. This covers $(i,j)=(2,1),(1,2)$ with their own coefficients. Their flow contributions are $-Q_1z^2/2-Qh^2/2+O(q\sqrt{E_J}+q^2)$. The $m$ and $n$ terms have coefficients $Q_1(1-m)/2$ and $Q(1-n)/2$; $|1-n|\le1+|n|$ gives the leading $QE_J$. Finally $|\langle P\Theta,N_{\rm sh}\rangle|\le\sqrt{2E_J}\,Cq\sqrt S(1+S)$ by REC's data-space shape projection. Adding the four pairings and this error proves \eqref{an04e5rv1:defect}.

Put $Y=(E_J+q^2)^{1/2}$. Division gives $Y'\le[Q/2+C(|n|+|1-m|)]Y+Cq$. From $t_b=t_h+b$ onwards variation of constants, $\int_0^\infty e^{-Qu/2}du=2/Q$, and the integrated budget give
\[
 \sup_{[t_b,t_f]}Y\le C_K e^{Q(t_f-t_b)/2}(Y(t_b)+q)
 \le C_K e^{-Lb}\big(1+\sqrt{E_J(t_b)}/q\big).
\]
Here $b\ge b_0$ is fixed before $q$; $t_f-t_h=\ell/L+O_K(1)$ and $q^2\ell\to0$. The last constant is independent of $b,J_2$.

We compute the entry on arbitrary subsequences at this \emph{fixed} $b$. HBD \texttt{clock} and BRW \texttt{field, delay} transfer compact right paths with mass $\mu_b=M(b)\pi$, $d\pi=4\cos^2\alpha\,dl$. The integrand of the right part of $E_J/q^2$ on a compact right interior is the physical radial mass times
$[(\sin(qx)/q)^2+(\sin(qy)/q)^2]/2$. It converges in averaged label law to $(x^2+y^2)/2$ integrated against $\mu_b$. At this same $b$ the all-label vertical bound $C_b q$ charges a strip of length $\sigma$ at most $C_b^2\sigma$ in scaled energy. Take its limsup at fixed $\sigma,b$, then remove this auxiliary strip at fixed $b$. The identical argument removes the Q2 and Q3 endpoint strips in the next calculation. On compact Q3 the squared radius is $O_{b,\sigma}(S_0)$ by the proof of COD \texttt{trained}; $S_0/q^2\to0$. On $J_2$, TBW \texttt{smallhit} gives $U_1,W_1=o_{b,J_2}(q)$. On compact Q2 the positive second components, balance and TBW \texttt{seed} give
$U_2/q,W_2/q\to c^\omega(4e^b)^L\sin\alpha$. Thus the Q2 contribution outside $J_2$ divided by $q^2$ tends to $c^{2\omega}(4e^b)^{2L}\int_{{\rm Q2}\setminus J_2}\sin^2\alpha\dd$.

It remains to bound the right moment after discount, uniformly over resident extractions. COD \texttt{ceiling} bounds $e^{-2Lb}(x_+^2+y_+^2)$ by $C(1+f^{-1/3})^2 e^{-\lambda fb/8}$ at normalized upper Q1 rank $f$, with $|df|=d\pi$, and with $f=1/2$ on Q4. The envelope is integrable at zero, so its integral tends to zero. For the negative parts COD's lower face in the proof of \texttt{exposure} gives $x_-\le C\sqrt{1+b}$ by order. The bottom $y$ equation in COD \texttt{field}, its bounded entry, its damping kernel $\le1$, and BRW \texttt{budget} give $(y_a)_-\le C+C\int_0^b e^{Lu}(1+u)^{-2/3}du$. After multiplication by $e^{-Lb}$, splitting at $b/2$ bounds this by $O(e^{-Lb/2}+(1+b)^{-2/3})$. Order bounds every $y_-$ by $(y_a)_-$. These bounds hold for every extracted law so also bound the original limsup. For uniformity, choose a sequence of parameters approaching the supremum at each fixed $b,J_2$, extract a limiting ratio, and use Lemma~\ref{an04e5rv1:matching}. The discount and the resident coefficient are evaluated at that same limiting ratio since $b$ is fixed. The resulting majorants use $L_-,\lambda_-$ and are independent of the limit point; only then send $b$ to infinity. It follows that
\[
 \limsup_{b\to\infty}\limsup_{q\downarrow0}
 \sup_{[t_b,t_f]}\sqrt{E_J}
 \le C_K\left(\int_{{\rm Q2}\setminus J_2}\sin^2\alpha\dd\right)^{1/2}.
\]
Exhaust Q2 only now. On this late interval the right vertical norm is at most $2\sqrt{E_J}$; the two cross responses have the bound derived above. For $T_1$ use $W_1$ on $J_2$, $W_2$ on its complement and $|v_1|\le C$. On $[t_0,t_b]$, for each fixed $b$ the all-label vertical envelope is $C_bq$, the mass bounded, $b_2=(U_2)_++O(q|U|)$ and $\|b_1\|\le C$. Thus the same quantities are $O_b(q)$ there. Given any tolerance choose the arc and fixed delay according to the iterated bound, then take $q$ small on both intervals. This proves uniform convergence on their union.
\end{proof}


For the strong response, the scalar four-regression equation has coefficient $a_m=Q_1\|\nabla m\|^2/2\ge Q_1m^2/S$. It is bounded below on the whole interval: before $\tau$ use Lemma~\ref{an04e5rv1:conelemma}, afterwards $Q_1(1-m)^2\le Q-\Delta$. The forcing in $(1-m)'+a_m(1-m)$ is $O(|z|+|h|+q)$ since $\langle g_m,g_n\rangle=O(qS)$ and the other gradients are bounded. It is uniformly $o(1)$ by Lemma~\ref{an04e5rv1:axislemma}. The entrance homogeneous term decays exponentially over the diverging time to every fixed half-centered window. Thus $m\to1$ there.

For the weak gradient fix such a window, with $K$ enlarged to include its endpoints and the bounded post-$\tau$ continuation. Put $\epsilon_J=\sup_{[t_b,t_f]}\sqrt{E_J}$. On $J_2$ up to $\tau$, the cone gives $U_2,W_2>0$ and balance gives
$|W_2-U_2|^2\le |W_2^2-U_2^2|\le W_1^2+U_1^2$.
After $\tau$, integrate the equations for $(U_2-W_2)/2$ and $(U_2+W_2)/2$ with coefficients $\mp\mathsf R_{22}$ and horizontal forcing bounded by $C(|W_1|+|U_1|)$. On this bounded interval both propagators are positive and bounded by a common constant. The difference mode in label norm, and the negative part of the sum mode, are therefore bounded by $C_K\epsilon_J$ on $J_2$, using the nonnegative sum at $\tau$. This bounds $\|\mathbf1_{J_2}(W_2-U_2)\|$ and $\|\mathbf1_{J_2}(U_2)_-\|$ there. On $J_2^c$ both vertical norms cost at most $2\epsilon_J$. Take the iterated limits of the lemma to obtain $\|W_2-(U_2)_+\|=o(1)$ on the centered window. Consequently $\|W_2-b_2\|=o(1)$. Also $|b_2(1-|v_2|^2)|\le Cq|U|$: use $|b_2-(U_2)_+|\le Cq|U|$, $|(v_2)_1|\le Cq$, and for positive $U_2$ bound $1-\theta_2$ by the weighted Gaussian wrong-gate tail $Ce^{-c(U_2/(q|U|))^2}$. Hence
$\int b_2^2|v_2|^2\dd=\int b_2^2\dd+O(qS)$. In the exact formula
$\|\nabla n\|^2=\int(b_2^2+W_2^2|v_2|^2)\dd$ replace $W_2$ by $b_2$ using the bounded norms, and use $n=\int W_2 b_2\dd$. This gives $\|\nabla n\|^2=2n+o(1)$. In CFCR's exact flow the $g_m$ pairing costs $O(qS)$, the $g_z,g_h$ terms vanish with $z,h$, and the shape remainder costs $O(q)$. Thus $n'=Q n(1-n)+o(1)$ uniformly on the window. Using the measured centering $n(t_{1/2})=1/2$ gives uniformly for bounded $s$
\begin{equation}\label{an04e5rv1:logistic}
 (m,z,h,n)(t_{1/2}+s)\longrightarrow(1,0,0,N(s)),\qquad
 N(s)=(1+e^{-\lambda s})^{-1}.
\end{equation}
\paragraph{Identification of the witness clock and the observable.}
WIT \texttt{orbit, clockalignment} defines a coherent weak \emph{mass} $n_\rho(s)$ on its own half-mass-centered clock. On the coherent leading state its value is also the weak regression: at the weak cluster center $D=\rho e_2$, the limiting strong direction $e_1$ gives zero, and unit weak directions $e_2$ give $f_2(\rho e_2)=\rho n_\rho$. Dividing $D_2f_2$ by $\rho^2$ gives $n_\rho$. Equivalently WIT \texttt{staticbounds} bounds this difference by $25qR$ for the compact coherent state on $J$. We apply no such chart bound to the trained trajectory. Its exact observable is $n_q=\E_2[D_2 f_2]/Q$, with $n_q(t_{1/2})=1/2$ by first hit and continuity. The preceding initialized argument, in the shifted variable $s=t-t_{1/2}$, gives the limiting ODE $N'=\rho^2 N(1-N)$ with $N(0)=1/2$. WIT has the same rate and initial value. Integrating the log odds (or using uniqueness) gives $N_\rho(s)=n_\rho(s)=(1+e^{-\rho^2 s})^{-1}$. The relative shift of these two limiting response profiles is zero. At finite noise the measured time translation is exactly $t\mapsto t-t_{1/2}$; the two finite-noise mass and regression observables need not coincide. Uniformity in \eqref{an04e5rv1:logistic} also follows by Gronwall on each fixed centered compact using $\lambda_+$ and the uniform $o(1)$ forcing.

By WIT \texttt{windows, rectangle, uniformmargins}, $J=[s_*-r,s_*+r]$ with a common $r>0$ for all of $\mathcal P_*$. Hence the chosen $s_*$ belongs to $\operatorname{int}J$ throughout, and
\begin{equation}\label{an04e5rv1:witness}
 N_\rho(s_*)\ge\tfrac14+\tfrac\sigma2,\qquad
 0\le(1-N_\rho(s_*))^2\le(\tfrac34-\tfrac\sigma2)^2
 <\tfrac9{16}=\tfrac{441}{784}<\tfrac{676}{784}=\tfrac{169}{196}.
\end{equation}
Here $0<\sigma=n_c(s_*-r)-1/4<3/4$. Cluster projection gives loss ratios $o(1)$ and $(1-N_\rho(s_*))^2+o(1)$, uniformly since $\rho^2\ge\lambda_-$ and $f_0(D)=S_0D/4$. Together with $m\to1$, this proves all four guards. Through $T$ a common lower bound for $1-n$ follows from $1-\eta$ before the hit and, on its bounded continuation, half the minimum of $1-N_\rho(s)$ on the fixed centered compact. For $|s|\le K_*$ the latter can be taken $[2(1+e^{\lambda_+K_*})]^{-1}$.

For each parameter define $e_*(\rho,\zeta)=\min_{s\in J}\mathcal E_{\rm coh}(\rho,\zeta;s)$. WIT's gate margin keeps the coherent formula active on $J$; \texttt{branchsmooth} gives joint continuity on $\mathcal P_*\times J$. The minimum is attained and its absolute value is bounded uniformly (for instance use WIT's $R$ and $\tfrac12 z_+^2+\tfrac12 R^2+R(z_++R)$ in \texttt{error}). The strict derivative bounds $\mathcal E'\le-3\gamma_0/4$ on $J_-$ and $\mathcal E'\ge3\gamma_0/4$ on $J_+$ exclude both endpoints as minimizers. Every minimizer lies in the interior. WIT fixes the common evaluation point $s_*$; the derivative signs place an interior minimizer for each parameter but do not identify it with $s_*$ away from the center.

\section{Probe exposure and the signed history to be estimated}\label{an04e5rv1:history}
Set $p_\alpha=(U_\alpha(T)^\top\xi)_+$, $F_i=f_{T,i}(\xi)$, $H=F_1/q$ and $w=Q(1-n)$. Before $t_d$ on the stopped interval, \eqref{an04e5rv1:envelope} bounds $\int\sqrt E\dt$. After that interval ends there are at most $O(1)$ units left to $T$, by the clock comparison. Consequently
\[
 \int_{t_0}^T(\|\Theta'\|+\|r\|_{P_1})\dt\le C,\qquad
 \|\mathcal U\|\le C,\quad
 \mathcal U_\alpha=\sup_{[t_0,T]}(|U_\alpha|+|W_\alpha|).
\]
Here the regression terms in $\|\Theta'\|$ are bounded by $C(|1-m|+|h|+|z|+\sqrt E+q)$ since $\|\nabla n\|\le C(\sqrt E+q)$. The strong residual costs $C(|1-m|+|z|+q)$ by shape projection. Minkowski's inequality bounds the time supremum in label norm.
On an excursion below $U_2=-A_0q\sqrt\ell\,\mathcal U_\alpha$, the $P_2$ gate has Gaussian upper tail and costs at most $Cq^2\mathcal U_\alpha$ per time in $U_2'$, for a sufficiently large fixed $A_0$. The $P_1$ cost is at most $Cq\mathcal U_\alpha\|r\|_{P_1}$ by Cauchy--Schwarz. Integrate from the last crossing or $t_0$, using the $O_{b_0}(q)$ entrance. Uniformly up to $T$,
\[
 (U_2)_-\le C_{b_0}q+Cq\sqrt\ell\,\mathcal U_\alpha,\qquad
 \|p\|\le Cq\sqrt\ell,\qquad |F_1|+|F_2|\le Cq\sqrt\ell .
\]
Define the strong-minus-weak history and the signed budget by
\begin{align}
 G_\alpha&=\int_0^T (Q_1 z b_1-w b_2)\dt,\nonumber\\
 M_q&=\frac{\int_{\rm Q4}pG_\alpha\dd-\int_{{\rm Q4}^c}p(-G_\alpha)_+\dd}{2q^2}.
 \label{an04e5rv1:budget}
\end{align}
The sign follows from $W_2'=(w b_2-Q_1zb_1)/2+\varepsilon_\alpha$. Orthogonal projection of both the residual and feature off $D_j$ in each cluster bounds $|\varepsilon_\alpha|\le Cq^2(1+S)|U_\alpha|$. Bounded mass through $T=O(\ell)$ therefore gives
$\int p|W_2(T)+G_\alpha/2|\dd\le C\|p\|(\sqrt{S_0}+q^2\ell)=o(q^2)$.
In particular $-F_2/q^2\ge M_q-o(1)$. The exact scaled loss is
\[
 e_q=H^2/2-\frac{\sin(q\zeta)}q H+
 \left(\cos(q\zeta)+F_2/2\right)F_2/q^2 .
\]
The signed endpoint criterion is
\begin{equation}\label{an04e5rv1:criterion}
 \liminf_{q\downarrow0}\inf_{\mathcal P_*}\left[M_q-H^2/2+\frac{\sin(q\zeta)}qH+e_*\right]>0 .
\end{equation}
For this implication, $A_q=H^2/2-q^{-1}\sin(q\zeta)H-e_*=O(1+\ell)$ and the vertical coefficient is $1+O(q\sqrt\ell)$. On the branch $F_2=o(1)$ the threshold for $-F_2/q^2$ is $A_q+o(1)$ uniformly: $|\sin(q\zeta)/q|\le z_+$, the coherent minimum is uniformly bounded, and multiplying $O(1+\ell)$ by the coefficient error $O(q\sqrt\ell)$ costs $o(1)$. The estimates below will give a diverging margin in \eqref{an04e5rv1:criterion}.

For the remainder set $A_R=\{\cos\alpha>0\}$, $A_L=\{\cos\alpha<0\}$,
\[
 C(t)=\int (U_1)_+(U_2)_+\dd,\quad I(t)=\int_{t_0}^t C(u)\,du,\quad
 B(t)=q(1+t-t_0),\qquad t_R=t_h+\ell/L-R .
\]
Here $C(t)$ denotes the overlap and unindexed constants $C$ in inequalities denote bounds. For all sufficiently large fixed $R$, $t_R$ precedes both $t_d$ and $\tau$ by the two-sided clock estimate. Equation~\eqref{an04e5rv1:envelope} then gives $\|U_2\|+\|W_2\|+\|b_2\|\le r(t)=C_0q e^{Q(t-t_0)/2}$ up to $t_R$, with $r(t_R)\le C_1e^{-LR}$ independently of fixed $R$ for small $q$.

\section{Uniform prefix estimates}\label{an04e5rv1:prefixsec}
The following details make the comparisons quantitative before the
terminal clock. Set $d_i=W_i-b_i$ and $w=Q(1-n)$. In the characteristic
equations $U'=\mathsf R^\top W,\ W'=\mathsf R U$, the regression
decomposition, including its uniform per-receiver shape error, gives
\[
 \mathsf R=\tfrac12 Q_1(1-m,-z)^\top v_1^\top+
            \tfrac12 Q(-h,1-n)^\top v_2^\top+O(q).
\]
Here the pairs denote column vectors in the outer products. The error
follows by data Cauchy--Schwarz against $D$ times the receiver gate and
the $O(q(1+S))$ cluster shape residual. Equivalently the output-layer
equation uses $b_i$ with error $O(q^2|U|)$ by orthogonality also of the
feature remainder. We will often only need $O(q|U|)$.
In particular $\mathsf R_{22}\ge-Cq$ and
$|\mathsf R_{12}|+|\mathsf R_{21}|\le C(|z|+|h|+q)$ uniformly.

For $x=1-m$, the scalar damped equations from the four-regression flow are
\[
 x'=-a_m x+O(|T_1 z|+|h|C+q),\qquad
 z'=-a_z z+\tfrac12 w C+\tfrac12 Q_1T_1 x+O(q).
\]
Here $a_z=Q_1\|\nabla z\|^2/2\ge Q_1\|b_1\|^2/2\ge Q_1m^2/(2S)$ by $m=\langle W_1,b_1\rangle$. Both coefficients have positive upper and lower bounds. On the incoming
interval $m>99/100$ and $S\le4$, so
\begin{equation}\label{an04e5rv1:damping}
 a_m\ge Q_1m^2/S>9801/40000,\qquad L\le49/200=9800/40000.
\end{equation}
Take $\eps=1/80000$ and $q^2/2<1/80000$. For sufficiently small $q$ the gap to $Q/2$ is bounded below, for example by $1/80000$. On the rest of the interval both damping coefficients stay positive since $m\ge1-\sqrt{Q/Q_1}\ge1/4$ after the small hit, and $S$ has a common bound.
The integrated cost of $T_1$ times the homogeneous entrance deficit
is $O_{b_0}(q)$: $|T_1|\le Cr(t)$ on the incoming interval, and the final
fixed interval starts after the entrance deficit has decayed to $o(q)$.
The coefficient comparison can use one \emph{fixed} incoming cutoff
before $R$ is chosen. Integrating the $z$ propagator up to an arbitrary
$t$ gives weight comparable to $1\wedge(t-u)$. Duhamel in the $x$
equation can be made explicit. Write $x=x_0 e^{-\int_{t_0}^t a_m}+x_f$, $x_0=x(t_0)$, and let $\beta=\sup(|T_1|+|h|)=o(1)$. Integrating the decaying kernel in $x_f$ gives
$\int|x_f|\le C\{\beta\int |z|+\beta I+q(t-t_0)\}$ on each prefix. Hence $\int|T_1 x|\le CB+o(1)(I+\int|z|)$. In the $z$ equation the $wC$ source is nonnegative, $w$ is bounded above and below, and $z(t_0)=O(q)$. Integrate first in the receiving time; the absolute forcing error is bounded by $C\int|T_1x|+CB$ in this integral norm. Absorb the $o(1)\int |z|$ into the absolute estimate, then substitute into the negative-part and signed estimates. Thus
\begin{align}
 \int_{t_0}^t |z|&\le C(I+B),&
 \int_{t_0}^t z_-&\le o_q(1) I+CB, \nonumber\\
 \int_{t_0}^t z&\ge c\int_{t_0}^t[1\wedge(t-u)]C(u)\,du-o_q(1)I-CB.
 \label{an04e5rv1:prefix}
\end{align}
Constants here and in the next estimates are independent of the compact
right-label truncation and the large incoming offset $R$.

Differentiating $b_i$ uses $v_i$, with $|v_1^\top v_2|\le Cq$,
$b_i=(U_i)_++O(q|U|)$ and
$|b_i(1-|v_i|^2)|\le Cq|U|$. The last bound uses
$\sup_{x>0}x\Phi(-x)<\infty$ in the positive coordinate.
The $d_1$ equation contains the terms
\[
 2d_1'=-Q_1x|v_1|^2d_1-Qb_2\langle b_2,d_1\rangle+
 Q_1 z W_2 |v_1|^2-Qb_2\langle b_1,b_2\rangle+O_{L^2}(q).
\]
At $t_0$, $\norm{d_1}=O_{b_0}(q)$: integrate the same inequality from
initialization, with $\norm{d_1(0)}=O(\sqrt{S_0})$, the HBD bounds for
the cross regressions and overlap $O(H_0V_0+qH_0^2)$, and the integrated
early shape cost $O(q\int H_0)=O_{b_0}(q)$, with the entrance envelopes of Section~\ref{an04e5rv1:setup}. The growth from $x_-$ is bounded by the early mass integral.
Dropping the dissipative rank-one term and using $\int |x|\le C$ gives
\[
 \norm{d_1(t)}\le C(I+B),\qquad
 \norm{d_1(t)}\le C(r(t)I+B)\quad(t\le t_R).
\]
Since $h=\langle b_1,b_2\rangle+\langle d_1,b_2\rangle$,
$\langle b_1,b_2\rangle=C+O(q)$ and $\int\norm{b_2}\le C$, we get
\begin{equation}
 \int_{t_0}^t|h|\le C(I+B),\qquad
 h_-(t)\le C r(t)^2 I(t)+C B(t)\quad(t\le t_R).
 \label{an04e5rv1:earlyh}
\end{equation}
Similarly $2d_2'=-w|v_2|^2d_2-Q_1zb_1+QhW_1|v_2|^2+O_{L^2}(q)$, so
$\norm{d_2}\le C(I+B)$. This dissipative bound will also control costly
returns outside the retained bulk.

For clarity, the terminal weight in \eqref{an04e5rv1:prefix} can be removed
up to constants and $O(B)$. Differentiation of
$C=\int(U_1)_+(U_2)_+\dd$ a.e. gives $|C'|\le C_2(C+I+B)$:
replace $W_i$ in each diagonal characteristic term by $b_i+d_i$,
use the two norm bounds, and use
$|z|+|h|\le C_2(C+I+B)$ for the off-diagonals. For this pointwise bound use the $z$ propagator bounded by one, $\int C=I$ and the just-established prefix bound for $\int|T_1 x|$. For $h$ use $h=\langle b_1,b_2\rangle+\langle d_1,b_2\rangle$.
With $I'=C$, Gronwall on an interval of length at most two,
starting at a point of $[t-2,t-1]$ where $C\le I(t-1)$, gives
$I(t)+C(t)\le C_3(I(t-1)+B(t))$ for $s\ge2$. Hence the weighted integral
is at least $cI(t)-CB(t)$.

\begin{lemma}[Two shrinking Q1 bands in the physical flow]\label{an04e5rv1:bandlemma}
Let $R$ be any sufficiently large fixed offset. Uniformly after a late fixed entrance,
\begin{equation}\label{an04e5rv1:uniform}
 \lim_{b\to\infty}\limsup_{q\downarrow0}\sup_{t_h+b\le t\le T} B(t)/I(t)=0 ,
 \qquad I(t_R-1)/q^{3/4}\longrightarrow\infty .
\end{equation}
\end{lemma}
\begin{proof}
We prove both assertions by testing the corresponding contrary sequences. First take $s=t-t_0\to\infty$ with $t\le t_R$ and $I(t)\le A B(t)$ for some fixed $A$. Use $\eps=(1+s)^{-4}$. In the second test take $t=t_R-1$, $I(t)\le A q^{3/4}$ and $\eps=q^{1/3}$. For either test write
\[
 g(u)=\sup_\alpha(|\mathsf R_{12}|+|\mathsf R_{21}|),\qquad
 P_t=\int_{t_0}^t g(u)\,du .
\]
By \eqref{an04e5rv1:prefix} and \eqref{an04e5rv1:earlyh}, $P_t=O_A(q(1+s))$ in the first test and $O_A(q^{3/4})$ in the second. In both,
\begin{equation}\label{an04e5rv1:units}
 q/\eps\to0,\quad P_t/\eps\to0,\quad
 \eps P_t/q\to0,\quad q e^{Qs/2}=O(1),\quad s=O(\ell).
\end{equation}
In particular $s\le\ell/L_-+C$. The band powers $4,1/3,3/4$ are absolute. The weak and strong units satisfy $qH/V\le C_{c,b_0}$ since $e^{-\omega t_0}/q=(c/(S_0e^{t_0}))^\omega$ stays between positive constants. The pre-entrance integrals use $1/\omega\le1/\omega_-$ and $1/(2L)\le1/(2L_-)$, and Gaussian costs use $\rho_-$ as the lower weak mean. In the second test $q e^{Qs/2}$ on the final unit is bounded below by $e^{-L_+(R+|b_0|+2)}$ for small $q$, using the exact cancellation $q e^{L\ell/L}=1$. In the first test $e^{Ls}/(1+s)^9\ge e^{L_-s}/(1+s)^9$. These estimates apply also to all parameter sequences in Lemma~\ref{an04e5rv1:matching}. For the first test use $q(1+s)^5\to0$ and $(1+s)^{-3}\to0$; for the second the two perturbation ratios have powers $3/4-1/3$ and $1/3+3/4-1=1/12$. All comparisons below are simultaneous for $\eps\le k=\cos\alpha\le2\eps$ on Q1.

For coordinate $i$ set $X_i=(U_i,W_i)$, $J(x,y)=(y,x)$, and $a_i=\mathsf R_{ii}$. Exactly,
$X_i'=a_i JX_i+(\mathsf R_{j i}W_j,\mathsf R_{i j}U_j)$ with $j=3-i$.
Its homogeneous matrix is $e^{J\int a_i}$, of max-norm $e^{|\int a_i|}$. On equal positive entrance components it multiplies each by $e^{\int a_i}$. In sum/difference coordinates $(p,d)=((U_i+W_i)/2,(U_i-W_i)/2)$ the multipliers are $e^{\int a_i}$ and $e^{-\int a_i}$. We estimate the exact inhomogeneous integrals on whole prefixes; strict improvements on a closed trial region below exclude a first contact at a single face or a corner.

\paragraph{Initialization to the fixed entrance.}
Write $a_1^*=1/2+q^2,\ a_2^*=L+q^2$, $H=H_0,\ V=V_0$ as in Section~\ref{an04e5rv1:setup}; $V/H=e^{-\omega u}$. At $t_0$, $H\asymp_{b_0}1$ and $e^{-\omega t_0}\asymp_{b_0}q$ since $S_0=cq^{1/\omega}$. Up to $t_0$, $qH\le C_{b_0}V$, $\int H^2\le C_{b_0}$ and $\int V^2\le C_{b_0}q^2$. HBD's physical argument gives $|X_1|\le C H$, $|X_2|\le C(V+qH)\le C V$ and $g\le C(q+HV+qH^2)\le C(q+HV)$ in max-norm. For the fixed extension beyond HBD's cap $8$, one can use the same scalar inequalities there: the spectral bound $S\le H^2$, the individual radial bound integrated with $\int H^2\le C_{b_0}$, and $D^+\sup|X_2|\le(a_2^*+CH^2)\sup|X_2|+CqH(1+H^2)$ give the stated bounds by integration. Absolute fitted diagonals cost $C S$, since $|f(D)|\le S|D|$. The target diagonal is nonnegative and at most $a_i^*=\E D_i^2$, hence $|a_i|\le a_i^*+CS$. All constants in this paragraph may depend on $b_0$.

Variation of constants in horizontal units $H$ charges the perturbation on any prefix by at most
\[
 C\int_0^{t_0}(V/H)(q+HV)\,du\le C(q/\omega+\int V^2)=O(q).
\]
It follows first that $|X_1|\le C\eps H$ without a sign assumption. In vertical units $V$ the perturbation now costs at most
\[
 C\eps\int_0^{t_0}(H/V)(q+HV)\,du
 \le C\eps\{q e^{\omega t_0}/\omega+\int H^2\}=O(\eps).
\]
These bounds include the factor $\exp(C\int S)$ for both propagators.

For lower bounds subtract each missed target diagonal from $a_i^*$. For positive $U_i$ the missed diagonal on the principal cluster is bounded by $C\exp[-c(U_i/(q|U|))^2]$, by Cauchy--Schwarz, the Gaussian half-space probability, and the bounded fourth data moment. The other cluster costs at most $q^2/2$. Thus
$a_i\ge a_i^*-CS-q^2/2-C\exp[-c(U_i/(q|U|))^2]$.
Choose $M_0$ larger than the common bound for $C\int_0^{t_0}S+1$, before choosing any lower faces. Use the two lower faces for each layer $X_1\ge a\eps H$, $X_2\ge a V$, componentwise, where $0<a<e^{-M_0-2}/8$. At initialization $k\ge\eps$ and $\sin\alpha\ge1/2$ for small $q$, so all faces are strict. On the stopped region the upper estimates give strong margin $\ge c_a\eps/q$, weak margin
$\ge c_a\min(1/q,e^{\omega(t_0-u)}/\eps)$. The first integrated tail is $O(\ell e^{-c_a\eps^2/q^2})=o(1)$ for either band. The second is at most
\[
 C\ell e^{-c_a/q^2}
 +C\int_0^\infty e^{-c_a e^{2\omega v}/\eps^2}\,dv=o(1);
\]
use $e^{2\omega v}\ge1+2\omega v$. This bound also covers delays $s$ growing arbitrarily slowly. The $q^2t_0$ cost vanishes. Accordingly $e^{\int_0^u a_i}$ is at least $e^{-M_0-1}e^{a_i^*u}$ on every stopped prefix for small $q$. Equal entrance components give individual lower bounds in horizontal and vertical units of $\eps e^{-M_0-1}-O(q)$ and $e^{-M_0-1}/2-O(\eps)$. Both strictly improve their trial bounds for both components, including simultaneous contacts. Boundedness and continuity continue the comparison up to $t_0$.

\paragraph{From $t_0$ to the test time.}
Set $v=u-t_0$, $V_*(v)=q e^{Qv/2}$. At $v=0$ we have $\underline a\eps\le U_1,W_1\le\bar a\eps$ and $\underline a q\le U_2,W_2\le\bar a q$ with $\underline a>0$. Balance gives $|U_1-W_1|\le Cq^2/\eps=o(\eps)$ there. The finite-noise regression field gives
\[
 |a_1|\le C(|1-m|+q),\qquad
 a_2=\tfrac Q2(1-n)\theta_2+O(q),\qquad a_2\ge-Cq .
\]
Use the global initialized budgets $\int|1-m|+\int |n|\le C$, $1-n>0$ and $q s=o(1)$. For any subinterval the horizontal propagator has norm bounded by $C$, and the vertical one by $C e^{Q(v-v')/2}$. To justify both upper faces without assuming them, put $X=\sup |X_1|/\eps$, $Y=\sup |X_2|/V_*$ over a prefix and the band. The two integral equations give
\[
 X\le C+C\,(\sup V_*) P_t Y/\eps,\qquad
 Y\le C+C\,\eps P_t X/q .
\]
Their product of perturbation coefficients tends to zero by \eqref{an04e5rv1:units}. Both suprema are bounded, and the differences from the homogeneous solutions on every prefix cost at most $O(P_t)=o(\eps)$ in absolute horizontal units and $O(\eps P_t/q)=o(1)$ in relative vertical units. These constants are independent of the labels in the test band.

The homogeneous horizontal sum component is at least $c\eps$ and the absolute difference component is $o(\eps)$ even after propagation. This proves individual horizontal lower bounds. For the two vertical lower faces use a trial bound $U_2,W_2\ge a' V_*$, where $a'>0$ will be chosen below. On this region $U_2/(q|U|)\ge c_{a'}\min(1/q,e^{Qv/2}/\eps)$. Thus
\[
 \int_0^s(1-\theta_2)\,dv
 \le C s e^{-c_{a'}/q^2}
     +C\int_0^\infty e^{-c_{a'}e^{Qv}/\eps^2}\,dv=o(1)
\]
by the same weighted Gaussian tail bound. On any prefix write $A_2=\int_0^v a_2(t_0+v')dv'$. Then
$A_2\ge Qv/2-M_1-o(1)$, $A_2\ge-Cqs$, where $M_1$ bounds $(Q/2)\int|n|$ independently of $a'$. If the initial vertical sum/difference modes are $p_0,d_0$, $p_0-|d_0|\ge\underline a q$, $|d_0|\le Cq$. Both homogeneous components are at least
\[
 e^{A_2}(p_0-e^{-2A_2}|d_0|)
 \ge e^{A_2}(\underline a q-o(q)).
\]
Choose $a'<\underline a e^{-M_1-2}/4$ before taking $q$ small. After the inhomogeneous error $o(V_*)$ both components exceed $2a'V_*$, initially and at every possible first contact. The upper bounds were obtained unconditionally; the strict margins cover simultaneous lower contacts and all labels with either band endpoint. This proves the comparison through the full stopped horizon.

\paragraph{The physical overlap and the two contradictions.}
The Q1 band has label measure at least $\eps/(2\pi)$ since $|d(\cos\alpha)/d\alpha|\le1$. On the last unit, $U_1\gtrsim\eps$ and $U_2\gtrsim q e^{Qs/2}$; the unit factor $e^{-Q/2}$ is bounded below. Therefore $I(t)\ge c\eps^2 q e^{Qs/2}$. In the first test this contradicts $I(t)\le A q(1+s)$ since $e^{Ls}/(1+s)^9\to\infty$. In the second, $q e^{Qs/2}$ at $t_R-1$ is bounded below by a positive constant for fixed $R,b_0$, so $I(t)\gtrsim_R q^{2/3}$ contradicts $I(t)\le A q^{3/4}$. This proves the second assertion and rules out every contrary first-test sequence before $t_R$. Beyond $t_R$, monotonicity of $I$ bounds $B/I$ by $O(q(1+\ell))/I(t_R-1)=o(1)$. Failure of the displayed iterated upper limit would yield $b_j\to\infty$, $q_j\to0$ and times with $s_j\to\infty$ and $I=O(B)$, which the two ranges exclude. No resident limits on shrinking label sets enter these comparisons.
\end{proof}
Together with \eqref{an04e5rv1:prefix} and the two-unit estimate, this gives $\int_{t_0}^t z\ge c I(t)$ and $\int_{t_0}^t z_-/I(t)$ as small as desired uniformly after a sufficiently late fixed delay.


\section{Choosing a bulk without losing exposure to discarded labels}\label{an04e5rv1:retained}
\begin{lemma}[Physical weight of retained and terminally active labels]\label{an04e5rv1:bulklemma}
For sufficiently small fixed $\delta>0$, remove open axis strips of total label measure at most $\delta$ from $A_R$, leaving two compact interiors. There exist a further exceptional subset of measure $o_{q;\delta}(1)$ and a retained set $K$ in these interiors such that, on $[t_0,T]$, with $k_\alpha=\cos\alpha$,
\begin{equation}\label{an04e5rv1:bulk}
 D^{-1}k_\alpha\le U_1,W_1,b_1\le D k_\alpha,\qquad |U|+|W|\le Dk_\alpha
 \quad(\alpha\in K).
\end{equation}
Here $D$ and the coefficient $C$ in $\|\mathbf1_{A_R\setminus K}\mathcal U\|\le C\sqrt\delta+o_{q;\delta}(1)$ are independent of $\delta,R,b$. Moreover $\|\mathbf1_{A_R}\sup_t(|U_2|+|W_2|)\|=o_q(1)$ and $\|\mathbf1_{A_L\cap\{p>0\}}\mathcal U\|=o_q(1)$.
\end{lemma}
\begin{proof}
Throughout let $g$ be the uniform off-diagonal bound used in Lemma~\ref{an04e5rv1:bandlemma} (equivalently its defining supremum). First $G=\int_{t_0}^Tg=o_q(1)$. To check the limit use an \emph{auxiliary} cutoff $t_{R'}$ at large fixed $R'$. On $[t_0,t_{R'}]$, $g\le C(r+q)$; integration costs $Ce^{-LR'}+O(q\ell)$. On the remaining interval of length at most $R'+C$, $g\le C(|z|+|h|+q)=o_q(1)$ by Lemma~\ref{an04e5rv1:axislemma}. Take $q\to0$ at fixed $R'$, then let $R'\to\infty$. The same argument with $\|\mathbf1_{A_R}X_2\|$ in place of $g$ gives $\int_{t_0}^T\|\mathbf1_{A_R}X_2\|dt=o(1)$, using the incoming energy envelope first and uniform right axis separation second. Neither limit fixes the $R$ to be used for the gate comparison.

Since $|\mathsf R_{22}|\le C$ and $\|\mathcal U\|\le C$, integrate the two vertical equations labelwise and apply Minkowski after taking the time supremum:
\[
 \|\mathbf1_{A_R}\sup_t |X_2(t)|\|
 \le \|\mathbf1_{A_R}|X_2(t_0)|\|
     +C\int_{t_0}^T\|\mathbf1_{A_R} X_2\|dt+C G\|\mathcal U\|=o(1).
\]
The first term is $O_{b_0}(q)$. Max- and sum-norms of the two components differ by at most a factor two. For the horizontal pair write $A_{1,\alpha}(u,v)=\int_u^v\mathsf R_{11,\alpha}$, with $|A_{1,\alpha}|\le M$ on all subintervals by the integrated $|1-m|+q$ bound. Hence the horizontal difference from its exact homogeneous evolution has time supremum $\le C G\mathcal U_\alpha$, a function of vanishing label norm. On each fixed compact right interior the entrance pair converges in $L^2(dl)$, uniformly over parameters, to $(a k_\alpha,a k_\alpha)$, $a=2\sqrt{M(b_0)}>0$, by HBD \texttt{clock} and BRW's fixed-delay continuation as matched in Lemma~\ref{an04e5rv1:matching}: the limiting radial weight is $4M(b_0)k_\alpha^2$, both angles tend to zero, and the integral of the squared error passes to the limit against the compact label-path law by boundedness. In particular, the maximum over $t$ of the error relative to the vector $e^{A_{1,\alpha}(t_0,t)}(a k_\alpha,a k_\alpha)$ has norm $o_{q;\delta}(1)$ there. The homogeneous multiplier is bounded between $e^{-M}$ and $e^M$ on this symmetric vector.

On the retained compact interiors discard labels where the preceding horizontal error or $\sup |X_2|$ exceeds $\epsilon_0 k_\alpha$, with $\epsilon_0<\min(1,a e^{-M}/4)$ fixed. By Chebyshev the discarded measure is $o_{q;\delta}(1)$ because $k_\alpha$ has a positive minimum on these \emph{fixed} compacts. On the good set both horizontal coordinates lie between $(a e^{-M}-\epsilon_0)k_\alpha$ and $(a e^M+\epsilon_0)k_\alpha$. The vertical upper bound and $|b_1-(U_1)_+|\le Cq|U|$ give all inequalities in \eqref{an04e5rv1:bulk} with one $D$ depending only on the global constants, $a,M,\epsilon_0$. No inverse minimum of $k_\alpha$ enters $D$. For any measurable $E\subset A_R$,
\[
 \|\mathbf1_E\mathcal U\|\le C\|\mathbf1_E X_1(t_0)\|+
 C G\|\mathcal U\|+C\|\mathbf1_{A_R}\sup_t |X_2(t)|\|
 \le C |E|^{1/2}+o_q(1)
\]
because all entrance horizontal coordinates have a common bound. Apply this to the axis strips and the additional discarded set. This proves small exceptional \emph{physical weight} in the norm of the time supremum, including its dependence on $\delta$.

On $A_L$, HBD's left mass bound gives $\|\mathbf1_{A_L}X_1(t_0)\|=O_{b_0}(q\ell)$. On any prefix the horizontal inhomogeneous solution costs at most $C\mathcal U_\alpha\int_{t_0}^t g$. Equations~\eqref{an04e5rv1:prefix} and \eqref{an04e5rv1:earlyh} therefore give
\begin{equation}\label{an04e5rv1:left}
 \|\mathbf1_{A_L}\sup_{[t_0,t]} |X_1|\|\le C(q\ell+I(t)).
\end{equation}
Over the full interval the same norm is $o(1)$ by $G=o(1)$. For the antisymmetric vertical mode $d=(U_2-W_2)/2$ the equation is $d'=-\mathsf R_{22}d+(\mathsf R_{12} W_1-\mathsf R_{21}U_1)/2$. Since $\mathsf R_{22}\ge-Cq$, its forward propagator is at most $e^{Cq(T-t_0)}$. Thus $\|\mathbf1_{A_L}\sup_t|d|\|\le C(q(1+\ell)+G\|\mathcal U\|)=o(1)$. The excursion bound in Section~\ref{an04e5rv1:history} gives $\|\sup_t (U_2)_-\|=O(q\sqrt\ell)$. Hence $d_2=W_2-b_2$ also has vanishing left norm uniformly in time, using the feature error $Cq|U|$.

On $D_L=A_L\cap\{p>0\}$, $U_2(T)<\tan(q\zeta) U_1(T)\le Cq|U(T)|$. Combining this upper bound with the negative excursion gives $\|\mathbf1_{D_L}U_2(T)\|=o(1)$. Therefore the terminal norm there of $v=(U_2+W_2)/2=U_2-d$ is $o(1)$. Its coefficient is $+\mathsf R_{22}$ and its forcing $(\mathsf R_{12}W_1+\mathsf R_{21}U_1)/2$. The backward propagator is bounded by $e^{Cq(T-t_0)}$. Integrating backward labelwise and then taking norms gives $\|\mathbf1_{D_L}\sup_t|v|\|\le C(\|\mathbf1_{D_L}v(T)\|+G\|\mathcal U\|)=o(1)$. Together with the estimates for $d$ and $X_1$ this yields
\begin{equation}\label{an04e5rv1:activeleft}
 \|\mathbf1_{D_L}\mathcal U\|=o(1).
\end{equation}
Every comparison for a retained right label uses $k_\alpha$ as unit. Left comparisons concern the integrated physical weight; the bounds hold including the negative-horizontal initial directions.
\end{proof}


\section{Entry into the negative-input region}\label{an04e5rv1:entry}
Write $\nu=\inf_{[t_0,T]}(1-n)>0$ with a common lower bound for small
$q$. For $\alpha\in K$, \eqref{an04e5rv1:bulk} and the prefix estimates give,
after a late fixed delay $b$,
\[
 \int_{t_0}^t z b_1\,du\ge c_1 k_\alpha I(t)
 \qquad(t_h+b\le t\le T).
\]
Indeed the positive part is multiplied by at least $k_\alpha/D$ and
the negative part by at most $D k_\alpha$. The constant $c_1$ works
independently of $\delta$. We next pay for weak excitation up to
$t_b=t_h+b$. At each fixed $b$ the scaled weak-feature integral on a
retained right compact satisfies
\[
 \limsup_{q\downarrow0}\sup_K
 \frac{1}{q k_\alpha}\int_{t_0}^{t_b}|b_2|\,du
 \le C_\delta(1+b)e^{(L-\nu_\delta)b},\qquad \nu_\delta>0.
\]
This follows by $|b_2|\le (U_2)_++Cq|U|$, fixed-interval tagged
transfer, and the COD rank ceiling for $x_+$ on ranks bounded away from
the upper endpoint; the Q4 bottom path obeys the same bound. Take
$\nu_\delta<L$ if necessary. In contrast BRW's protected positive
first moment on the last unit gives
$\liminf_{q\to0} I(t_b)/q\ge c e^{Lb}/(1+b)$ for all large fixed $b$.
To transfer this lower bound, retain the positive overlap from the
right labels and use fixed-time extraction:
$(U_1)_+(U_2)_+/q$ converges there to the radial mass times $x_+$;
the nonnegative remainder may be discarded. Use TBW \texttt{an04tbwv1:bulk} with $F_b=C_F\log(1+b)/(1+b)$. At fixed truncation the minimum normalized upper rank $f_\delta$ on compact Q1 is positive; Q4 uses $f=1/2$. Choose $b$ so $F_b<\min(f_\delta,1/2)$ and $b-B_F\log(1+b)>0$. All these receivers lie outside the exceptional upper ranks $f<F_b$. The limiting upper input is $\le-\sqrt{b-B_F\log(1+b)}<0$. For varying labels in a fixed interior compact the same rank inequality holds on each tagged limiting path by same-field order. A violating sequence at this fixed $b$ would have such a tagged limit, so all retained compacts have $U_2<0$ at $t_b$ uniformly for small $q$.
Thus, after fixing the compacts, choose $b$ sufficiently large for
the initial weak integral and $|W_2(t_0)|$ to cost less than any
prescribed fixed fraction of $k_\alpha I(t_b)$.

As long as $U_2\le0$ after this entrance on a retained label,
$|b_2|\le Cq k_\alpha$. The output equation integrates to
\begin{equation}
 W_2(t)\le -c_2 k_\alpha I(t),\qquad t_b\le t\le T,
 \label{an04e5rv1:wnegative}
\end{equation}
on this stopped interval: the additional weak cost is $Cq(t-t_b)k_\alpha$,
absorbed uniformly by \eqref{an04e5rv1:uniform}; the output shape error is
smaller. Explicitly, reserve $a_s=c_1/2$ from the lower bound for $(Q_1/2)\int z b_1/(k_\alpha I)$, since $Q_1\ge1$. Charge at most $a_s/8$ each to the entrance value of $W_2$, the weak term before $t_b$, the subsequent weak term, and the integrated shape error in these units. For the last two charges use respectively $CB/I$ and $CqB/I$ on each prefix. For the first two use the fixed-delay bounds above and $I(t)\ge I(t_b)$. The positive constant $c_2=a_s/2$ is fixed before $R,\delta,b$; only enforcement of the four fractions needs their subsequent choices. The input equation on that interval is
\begin{equation}
 U_2'=\tfrac Q2[(1-n)W_2-hW_1]\theta_2+O(qk_\alpha).
 \label{an04e5rv1:gate}
\end{equation}
Fix $\kappa>0$ sufficiently small compared with $\nu c_2/D$.
Whenever $h_-\le\kappa I$, \eqref{an04e5rv1:gate} gives
$U_2'\le-c_3 k_\alpha I\theta_2+Cq k_\alpha$.

For a sufficiently small fixed $d_g>0$ the stronger face needed at the final passage is
\begin{equation}\label{an04e5rv1:deepgate}
 U_2\le -d_g q k_\alpha\sqrt\ell .
\end{equation}
The next lemma proves entry by $t_R$ and continuation through $T$.

\section{Closing the final interval}\label{an04e5rv1:centered}
\begin{lemma}[Continuation of the negative input face]\label{an04e5rv1:gatelemma}
In the order of choices below, every label retained in Lemma~\ref{an04e5rv1:bulklemma} satisfies \eqref{an04e5rv1:deepgate} for $t_R\le t\le T$, and \eqref{an04e5rv1:wnegative} holds from $t_b$ through $T$.
\end{lemma}
\begin{proof}
The depth constant $d_g$ is distinct from the amplitude cutoff $d$ of Section~\ref{an04e5rv1:mass}. Fix $\nu>0$ from the passage bound $1-n\ge\nu$ and take $\kappa$ so that $\kappa D<\nu c_2/2$. Section~\ref{an04e5rv1:entry} gives the entry with these constants: on $U_2=0$ the principal weak gate is independent of $D_2$ and has probability $1/2$ (the retained $U_1>0$), so $\theta_2=1/2$. At $t_b$ the compact entrance lies strictly below zero. Up to $t_R$, \eqref{an04e5rv1:earlyh} gives $h_-/I\le C r(t_R)^2+CB/I<\kappa$ after $t_b$ with the choices below. While such an input is nonpositive, \eqref{an04e5rv1:wnegative} and this bound on $h_-$ give strict negative velocity at zero and on its interior side sufficiently near zero. One may also stop each individual label here because \eqref{an04e5rv1:earlyh} holds independently of other labels' gates. This gives $U_2\le0$ up to $t_R$. On $[t_R-1,t_R]$ put $u_\alpha=U_2/k_\alpha$. In the strip $-d_gq\sqrt\ell\le u_\alpha\le0$, condition on $|D_2-\rho|\le q$ in $P_2$. This has probability $p_0=\Pr(|Z|\le1)>0$, and $D_2>\rho/2$ for $q<\rho_-/2$. Given $D_2$, $D_1/q$ is an independent standard normal. Since $U_1\ge k_\alpha/D$, its gate threshold $-U_2D_2/(qU_1)$ is at most $D d_g(\rho+q)\sqrt\ell$. Choose $\gamma=1/8$, $d_g=1/(4\rho_+D)$, so $(\rho D d_g)^2/2\le1/32<\gamma$. Integrating the density on the unit interval beyond the upper threshold gives at least $(2\pi)^{-1/2}\exp[-(1+Dd_g(\rho+q)\sqrt\ell)^2/2]\ge c q^\gamma$ for small $q$, uniformly. On the slab $D_2^2/Q\ge(\rho_-/2)^2/(\lambda_++1/100)$, proving the same uniform lower bound for $\theta_2$ after multiplication by $p_0$. The normalized speed is $\le-c I(t_R-1)q^\gamma+Cq$. Since $I(t_R-1)\gg q^{3/4}$, the magnitude of this inward speed exceeds the full normalized strip width divided by one unit of time for small $q$. Labels starting inside it cross the deep face before $t_R$; labels already below cannot cross upward. The inequalities include the endpoints and are uniform on $K$.

To continue, stop at the first time the supremum on $K$ of $u_\alpha$ reaches $-3d_gq\sqrt\ell/4$, or at $T$. Initially the supremum is at most $-d_gq\sqrt\ell$. It is Lipschitz in time because the field is bounded and $|W|/k_\alpha\le D$ on $K$. This gives a genuine stopped interval even without attaining the supremum. On it all retained labels are strictly negative. Their weak gate satisfies $|v_2|^2=o_q(1)$ uniformly: the margin magnitude is at least $3\rho d_g\sqrt\ell/(4D)$, Cauchy--Schwarz bounds the weighted principal tail by a constant times the square root of its probability, and $|(v_2)_1|\le Cq$. In particular the comparison for $W_2$ applies on every prefix from $t_b$.

Let $d_2=W_2-b_2$. The exact definition $T_2=\int W_1W_2|v_2|^2\dd$ and the feature bound $|b_2(1-|v_2|^2)|\le Cq|U|$ give $T_2-h=\int W_1|v_2|^2d_2\dd+O(q)$. Split this integral three ways. On $K$ it costs $o(1)\|W_1\|\|d_2\|$. On $A_R\setminus K$ it costs $(C\sqrt\delta+o_{q;\delta}(1))\|d_2\|$. On $A_L$ it costs $C(q\ell+I)\|\mathbf1_{A_L}d_2\|=o(1)(q\ell+I)$. Here $\|d_2\|\le C(I+B)$ by Section~\ref{an04e5rv1:prefixsec}. Since $I(t)\ge I(t_R-1)\gg q^{3/4}$ and $B=O(q\ell)$, uniformly on every stopped prefix,
\begin{equation}\label{an04e5rv1:t2}
 |T_2-h|\le(C_0\sqrt\delta+o_{q;R,\delta,b}(1))I(t).
\end{equation}
$C_0$ is independent of $R,\delta,b$; only the threshold for small $q$ can depend on all three.

Now the scalar CFCR equation with its $O(q)$ remainder gives
\[
 h'=a_h h+\tfrac12 Q(1-n)(T_2-h)+\tfrac12 Q_1(1-m)C(t)+O(q),
 \qquad a_h=\tfrac Q2(1-n-\|\nabla h\|^2).
\]
For every fixed $R$, $[t_R,T]$ lies in a fixed half-centered compact; \eqref{an04e5rv1:logistic} gives $n>0$ and $1-m=o_{q;R}(1)$ there. Consequently $a_h\le Q/2$. Also $C(t)\le C'(I+B)$, $q/I=o(1)$. By \eqref{an04e5rv1:earlyh},
$h_-(t_R)\le(C_1 e^{-2LR}+o_{q;R}(1))I(t_R)$, $C_1$ independent of $R$. Write $T-t_R\le R+C_T$ with $C_T$ fixed by the clock. The scalar propagator is positive. Its maximum is $e^{Q(R+C_T)/2}\le 2e^{L(R+C_T)}$ for small $q$ at fixed $R$. Applying variation of constants, using monotonicity of $I$ for both terms, gives on each stopped prefix
\[
 h_-(t)/I(t)\le C_2e^{-LR}+C_R\sqrt\delta+o_{q;R,\delta,b}(1),
 \qquad C_R=C_3(1+R+|C_T|)e^{L_+(R+|C_T|)}.
\]
The constants $C_2,C_3$ are independent of the rectangle point, $R,\delta,b$. In the entry term combine $e^{-2LR}$ with $e^{LR}$ at the same $\rho$ before bounding by $e^{-L_-R}$. The displayed $C_R$ bounds the forcing at every rectangle point. This estimate was derived assuming only the \emph{outer} stopped face; it improves the required $h_-$ bound there.

\begin{center}\small
\begin{tabular}{p{.14\linewidth}p{.76\linewidth}}
Choice & Dependencies and inequalities\\ \hline
First & Fix the family, $b_0$, passage cutoffs, $C_T,D,\nu,c_2,\kappa,\gamma,d_g$. The coefficient $c_2$ comes from the signed prefix with global $D,\nu$ and fixed relative error fractions.\\
$R$ & Take large for the incoming envelope, $Cr(t_R)^2<\kappa/4$ in \eqref{an04e5rv1:earlyh} and $C_2e^{-LR}<\kappa/4$. This fixes $C_R$.\\
$\delta$ & Keep a fixed central Q4 subarc and require $C_R\sqrt\delta<\kappa/4$, $C_5\sqrt\delta<c_4/2$ with the endpoint constants of Section~\ref{an04e5rv1:endpoint}.\\
$b$ & At fixed compacts, take $b$ large for TBW's negative entrance and the COD/BRW weak-cost comparison. Require the uniform $B/I$ and negative-$z$ costs small for the chosen fractions, including $CB/I<\kappa/4$.\\
$q$ & At fixed $R,\delta,b$ enforce strict tagged entrance, $t_b<t_R-1$, the band lower bound, and all remaining error fractions, in particular $|o_{q;R,\delta,b}(1)|<\kappa/4$.\\
\end{tabular}
\end{center}
The endpoint constants used here can be fixed before $\delta$: choose a closed central Q4 arc of positive length with $k_\alpha$ bounded below. Its retained measure is at least half its length for small $q$ at fixed truncation. For example, fix the Q4 arc $-\pi/3\le\alpha\le-\pi/6$ modulo $2\pi$, of label measure $1/12$, with $k_\alpha\ge1/2$. Its retained measure is at least $1/24$ and the exposure bound of Section~\ref{an04e5rv1:endpoint} permits $c_4=d_gc_2/96$. $C_5$ uses the global $d_2$ and excursion bounds and the absolute physical strip constant of Lemma~\ref{an04e5rv1:bulklemma}. Neither requires an endpoint conclusion to set its value.

With these choices $h_-<\kappa I$ on the stopped interval. Equation~\eqref{an04e5rv1:gate} with $W_1\le Dk_\alpha$ gives $u_\alpha'\le-c I\theta_2+Cq$. Throughout the strip $-d_g q\sqrt\ell\le u_\alpha\le-3d_g q\sqrt\ell/4$ the same lower tail $\theta_2\ge c' q^\gamma$ applies, while the uniform upper-tail estimate already paid for \eqref{an04e5rv1:t2}. The derivative is strictly negative on the entire strip. The positive part $(u_\alpha+d_g q\sqrt\ell)_+$ starts at zero and cannot increase on the stopped interval, by its a.e. derivative, including both faces. Thus every retained label stays at or below the inner face. Taking the supremum leaves a positive distance to the outer stopping level even at a putative first exit. The stop must extend to $T$, including the case of exit and $T$ coinciding.
\end{proof}


\section{Subtracting all positive output outside the core}\label{an04e5rv1:endpoint}
For $qz_+\le1$, $\sin(q\zeta)\ge qz_-/2>0$ and $0\le1-\cos(q\zeta)\le q^2 z_+^2/2$. On $K$, $p\ge qk_\alpha z_-/(2D)+d_gq k_\alpha\sqrt\ell\,\cos(q\zeta)\ge d_g q k_\alpha\sqrt\ell$ since $d_g q^2 z_+^2\sqrt\ell\to0$. Thus the exact exposure has the same lower face constant, and
$-W_2(T)\ge c_2 k_\alpha I(T)$. The retained Q4 contains, for small
$\delta$ and small $q$, all but vanishing measure of a fixed central
Q4 subarc. Thus its favorable contribution is at least
$c_4 q\sqrt\ell\,I(T)$ with $c_4>0$ independent of the truncation.
Every retained Q1 label also has negative output.

To charge \emph{all} adverse labels, including any exceptional Q4,
let $D_*=\{p>0\}\setminus K$. Equations \eqref{an04e5rv1:activeleft} and
the right-strip estimate give
$\norm{\mathbf1_{D_*}\mathcal U}\le C\sqrt\delta+o_q(1)$.
The exposure excursion bound consequently gives
$\norm{\mathbf1_{D_*}p}\le(C\sqrt\delta+o_q(1))q\sqrt\ell$.
On the probe-active side $b_2\le Cq|U|$ because $U_2\le Cq|U|$.
Hence $\norm{\mathbf1_{\{p>0\}}(W_2)_+}\le\norm{d_2}+Cq\sqrt S
\le C I(T)$. Cauchy--Schwarz pays the entire adverse output by
$(C_5\sqrt\delta+o_q(1))q\sqrt\ell\,I(T)$.
Take $\delta$ small enough also that $C_5\sqrt\delta<c_4/2$.
The constants here precede the compact-specific choice of $b$.

Finally let $G_\alpha=\int_0^T[Q_1zb_1-w b_2]\dt$ and define
$M_q=(\int_{\rm Q4}pG_\alpha\dd-\int_{{\rm Q4}^c}p(-G_\alpha)_+\dd)/(2q^2)$.
The output projection error in replacing $-W_2(T)$ by $G_\alpha/2$
in these weighted integrals is $o(q^2)$ in absolute integral. Indeed
the shape error costs $O(q^2\ell)\norm p$, and the initial term costs
$O(\sqrt{S_0})\norm p$. Positive part is 1-Lipschitz. Consequently
\[
 M_q\ge c\sqrt\ell\, I(T)/q-o(1),\qquad
 H^2=O(\ell),\qquad I(T)/q^{3/4}\longrightarrow\infty .
\]
Since $\sqrt\ell\,I(T)/q\gg\ell$, this proves \eqref{an04e5rv1:criterion} with diverging margin. Directly, if
$F_i=f_{T,i}(\xi)$, then
$e_q=H^2/2-q^{-1}\sin(q\zeta)H+
(\cos(q\zeta)+F_2/2)F_2/q^2$.
The last coefficient tends to one by $\norm p=O(q\sqrt\ell)$.
For fixed $R$, the contradiction in \eqref{an04e5rv1:uniform} excludes every sequence with bounded $I(t_R-1)/q^{3/4}$ even if its parameters vary. Thus $\inf_{\mathcal P_*} I(T)/q^{3/4}\to\infty$. Together with the uniform upper bound $H^2=O(\ell)$ this gives $\sup_{\mathcal P_*}e_q(T)\to-\infty$, hence the strict comparison with
the finite coherent minimum at the measured learned endpoint. Since $F_i=O(q\sqrt\ell)$ the unscaled probe loss tends to $1/2$, as asserted in Theorem~\ref{an04e5rv1:main}.
\end{document}
